\documentclass[11pt,oneside]{article}
\usepackage[T1]{fontenc}
\usepackage[utf8]{inputenc}
\usepackage{lmodern}
\usepackage[margin=1in]{geometry}
\usepackage{microtype}
\usepackage{amsmath,amssymb,amsthm,mathtools}
\usepackage{booktabs,longtable,array,tabularx}
\usepackage{enumitem}
\usepackage{xcolor}
\usepackage{listings}
\usepackage{fancyhdr}
\usepackage{hyperref}
\usepackage[nameinlink,noabbrev]{cleveref}
\usepackage{url}
\usepackage{csquotes}
\usepackage{caption}
\usepackage{multicol}

\definecolor{linkblue}{RGB}{30,80,145}
\definecolor{codegray}{RGB}{245,245,245}
\hypersetup{
  colorlinks=true,
  linkcolor=linkblue,
  citecolor=linkblue,
  urlcolor=linkblue,
  pdftitle={Proof-Producing Finitary Interface Graphs for iPRA},
  pdfauthor={OpenAI research synthesis for the ProveIt project}
}
\setlist[itemize,enumerate]{nosep,leftmargin=*}
\setlist[description]{nosep,leftmargin=2em,labelindent=0pt,style=nextline,font=\normalfont\bfseries}
\setlength{\headheight}{14pt}
\setlength{\emergencystretch}{3em}
\setlength{\parindent}{0pt}
\setlength{\parskip}{0.58em}
\setcounter{tocdepth}{2}
\setcounter{secnumdepth}{3}
\pagestyle{fancy}
\fancyhf{}
\fancyhead[L]{PFIGs for iPRA}
\fancyhead[R]{\thepage}
\renewcommand{\headrulewidth}{0.4pt}

\lstset{
  basicstyle=\ttfamily\small,
  backgroundcolor=\color{codegray},
  frame=single,
  breaklines=true,
  columns=fullflexible,
  keepspaces=true,
  showstringspaces=false,
  keywordstyle=\bfseries,
  commentstyle=\itshape
}

\newtheorem{theorem}{Theorem}[section]
\newtheorem{lemma}[theorem]{Lemma}
\newtheorem{proposition}[theorem]{Proposition}
\newtheorem{corollary}[theorem]{Corollary}
\newtheorem{definition}[theorem]{Definition}
\newtheorem{principle}[theorem]{Principle}
\newtheorem{barrier}[theorem]{Barrier}
\newtheorem{example}[theorem]{Example}
\theoremstyle{remark}
\newtheorem{remark}[theorem]{Remark}
\newtheorem{status}[theorem]{Status}

\newcommand{\iPRA}{\mathrm{iPRA}}
\newcommand{\PRA}{\mathrm{PRA}}
\newcommand{\Nat}{\mathbb{N}}
\newcommand{\Bool}{\mathbb{B}}
\newcommand{\Prf}{\mathsf{Prf}}
\newcommand{\Good}{\mathsf{Good}}
\newcommand{\Edge}{\mathsf{Edge}}
\newcommand{\Viable}{\mathsf{Viable}}
\newcommand{\Obs}{\mathsf{Obs}}
\newcommand{\Pre}{\mathsf{Pre}}
\newcommand{\Post}{\mathsf{Post}}
\newcommand{\Demand}{\mathsf{Demand}}
\newcommand{\Code}{\mathsf{Code}}
\newcommand{\PFIG}{\mathsf{PFIG}}
\newcommand{\PRAAx}{\mathsf{PRAAx}}
\newcommand{\PTm}{\mathsf{PTm}}
\newcommand{\Fm}{\mathsf{Fm}}
\newcommand{\qf}{\mathrm{qf}}
\newcommand{\eps}{\varepsilon}
\newcommand{\del}{\delta}

\title{\textbf{Proof-Producing Finitary Interface Graphs}\\
A General Framework for Embedding Mathematics in Intuitionistic Primitive Recursive Arithmetic}
\author{Research synthesis prepared for the ProveIt iPRA development}
\date{15 November 2026\\Repository snapshot: \texttt{34816faca8aa605999290b41d26210753b5c3b38}}

\begin{document}
\maketitle

\begin{abstract}
The ProveIt repository contains a large collection of proposed reductions of algebra, analysis, number theory, geometry, probability, and topology to intuitionistic primitive recursive arithmetic (iPRA).  The recurring difficulty is not usually one isolated arithmetic identity.  It is the interface between ordinary mathematics and the very small object theory: approximate objects are represented by codes; computations must be primitive recursive; induction is restricted to quantifier-free formulas; non-unique witnesses must not be silently coerced into coherent choices; and a semantic Lean theorem is not yet an object-level derivation of type \(\Prf\ \PRAAx\ \Gamma\ \varphi\).

This report develops \emph{proof-producing finitary interface graphs} (PFIGs).  A PFIG is a fixed finite acyclic graph of exact transducers, bounded loops, bounded searches, relational producers, guarded partial operations, and stable observers.  Each node carries a backward precision-demand transformer and a local object-iPRA certificate.  The main normalization theorem turns every fixed PFIG into a query-specialized encoded state machine with one bounded primitive recursion and one quantifier-free invariant.  Consequently, the already available object-language recursion theorem can serve as the generic proof-producing backend; no internal universal primitive-recursive evaluator is required.

The framework gives complete human proofs of bounded Booleanization, least bounded search, trace reflection, demand propagation, relational viability, finite witness amalgamation, extraction of coherent regular names from shrinking viable clouds, stable-observer erasure, guarded branch stabilization, finite-cover gluing, and the equivalence between explicit height bounds and complete finite lists.  It also completes, at the mathematical level, the negative-translation argument needed to transfer retained classical PRA proofs into iPRA.  The report distinguishes genuine resolutions from principled reformulations and from open producer obligations such as general effective height bounds.  A detailed Lean development plan, source-oriented theorem signatures, exact-arithmetic Python reference implementation, regression tests, theorem ledger, and difficulty matrix accompany the paper.
\end{abstract}

\vspace{0.5em}
\noindent\textbf{Verification status.}
The mathematical arguments in this report are intended to be complete at the human level.  The Python reference tests included in the archive were executed successfully.  The Lean files are a source-oriented formalization blueprint and are not claimed to compile in this artifact environment.  Existing repository theorems explicitly cited as object-level results are inherited from the inspected snapshot; this packet does not claim to have rechecked the repository.

\tableofcontents
\newpage

\section{Scope, repository audit, and the problem to be solved}
\label{sec:scope}

\subsection{The target theory}

The repository uses \(\iPRA\) for first-order intuitionistic logic with equality, one function symbol for every primitive-recursive definition, the defining equations for those symbols, the usual successor properties, and induction restricted to quantifier-free formulas.  In the checked \texttt{ConciseIPRA} development, ordinary iPRA derivability is represented by
\[
   \Prf\ \PRAAx\ \Gamma\ \varphi,
\]
and is connected by a kernel-checked bi-interpretation to a much smaller native scalar calculus \cite{repo-pra-readme,repo-concise}.  This is a strong backend: ordinary mathematical arguments may be stated in the familiar \(\Prf\ \PRAAx\) layer and transported to the small native presentation only after a completed object derivation exists.

The restriction that matters most is not merely the absence of classical logic.  It is that the induction schema accepts only quantifier-free matrices.  Consequently, a proof that says ``for every precision there exists a sufficiently accurate approximation'' cannot be inserted unchanged as an induction invariant.  It must first be compiled into finite numerical state, primitive-recursive Boolean checks, and a quantifier-free invariant.  Likewise, semantic truth in Lean or a successful Python computation is not, by itself, a value of the inductive type \(\Prf\ \PRAAx\ \Gamma\ \varphi\).

\subsection{What is already present}

At the inspected snapshot, the directory contains two very different classes of material.

\begin{enumerate}
\item The \texttt{ConciseIPRA} subtree contains checked Lean definitions of syntax, natural deduction, primitive-recursive terms, iPRA axioms, native calculi, normalization, and mutual interpretations.  In particular, \texttt{Deduction.lean} provides substitution, weakening, cut, equality transport, and the ordinary logical rules \cite{repo-deduction}.
\item The \texttt{docs/proofs} subtree contains seventy-two delivered research packets across algebra, analysis, analytic and Diophantine number theory, geometry, probability, topology, combinatorics, and foundations.  The repository README explicitly records that these packets were filed as delivered and were not globally rerun or verified.  Most provide a mathematical reduction, reference code, and a Lean plan; only a small number construct actual \(\Prf\ \PRAAx\) values for framework combinators or toy examples, and none claims such a value for its main theorem \cite{repo-pra-readme}.
\end{enumerate}

Several foundation packets already identify the right ingredients: bounded Booleanization, proof-producing traces, finite observations, relational answer clouds, backward viability, coherent shrinking, precision contracts, stable observables, finite gluing, and explicit producer/verifier boundaries.  The most recent unified packet gives a broad synthesis, but its generic trace compiler and relational Lean files remain blueprints with visible gaps \cite{repo-unified}.  Other packets provide important checked or source-complete fragments.  Three are especially useful:

\begin{itemize}
\item \texttt{ObjectComposition.lean} defines exact term contracts and composes their actual object-iPRA derivations.
\item \texttt{ObjectIteration.lean} proves a generic primitive-recursion invariant theorem whose hypotheses and conclusion are genuine \(\Prf\ \PRAAx\) derivations \cite{repo-local-global}.
\item \texttt{ObjectObservation.lean} combines independently produced finite certificates through a local stability proof, again entirely at the object-derivation level \cite{repo-observable}.
\end{itemize}

These fragments suggest that the main missing bridge need not be a large new logical kernel.  It should instead be a disciplined compilation architecture that reduces diverse mathematics to a small collection of already-supported proof constructors.

\subsection{Recurring failure modes}

The repository reports are remarkably consistent about where reductions become unreliable.  The following list is not a list of unrelated defects; it is the design specification for a general framework.

\paragraph{Ambient proof versus object proof.}
A Lean theorem about natural numbers, lists, rationals, or finite graphs may establish the intended fact in the metatheory while never constructing the corresponding object formula or its \(\Prf\ \PRAAx\) derivation.  A trustworthy pipeline must make the change of level explicit.

\paragraph{Quantifiers inside induction.}
Ordinary correctness statements quantify over witnesses, precisions, partitions, roots, covers, or paths.  The iPRA induction predicate must instead be a quantifier-free relation on a finite encoded state.  Bounded quantifiers must be turned into primitive-recursive folds before induction is invoked.

\paragraph{A hidden universal evaluator.}
A proof plan may speak of ``the code of an arbitrary primitive-recursive program'' and then evaluate it by one total primitive-recursive function.  No such total universal evaluator exists.  Fixed families may be specialized externally, and clocked traces may be checked, but the framework must not smuggle in universality.

\paragraph{Non-unique witnesses forced into a single name.}
At each finite precision there may be many valid roots, eigenvectors, chart representatives, local solutions, or combinatorial certificates.  Choosing the least code independently at every precision generally destroys coherence.  A finite relation is legitimate; an infinite coherent branch requires additional quantitative hypotheses.

\paragraph{Discontinuous labels.}
Root orderings at collisions, eigenvector signs, rotation axes at the identity, exact sides of zero, and chart choices are not uniformly continuous.  They cannot be made total by notation.  One must use invariant observables, clusters, projective data, or an explicit separation guard.

\paragraph{Partial operations justified by mere negation.}
Division, normalization, pivot selection, logarithms, inverse charts, and spectral projectors require positive quantitative margins, not merely a proof of \(x\neq 0\).  A constructive interface must consume an apartness bound or return a relation of admissible outcomes.

\paragraph{Compactness and convergence without data.}
Abstract compactness, ``take a convergent subsequence,'' asymptotic notation with ineffective constants, and metastability do not provide primitive-recursive moduli.  Finite nets, covering certificates, tail bounds, and convergence rates must be present as inputs or separately produced certificates.

\paragraph{Finiteness confused with a search cutoff.}
A proof that a Diophantine set is finite does not yield a primitive-recursive height bound.  A cardinality bound does not locate the points.  The finite verifier and the height producer are different theorems.

\subsection{Design goals}

A useful general framework should satisfy all of the following.

\begin{enumerate}
\item \textbf{Object-level closure.}  Local correctness certificates have type \(\Prf\ \PRAAx\), and the generic compiler returns a term of the same type.
\item \textbf{Fixed specialization.}  The mathematical construction is a fixed finite graph.  Its primitive-recursive symbols are generated externally; no universal evaluator is postulated internally.
\item \textbf{Finite-use semantics.}  A terminal precision request causes a backward computation of exactly which finite observations are needed.
\item \textbf{Relational honesty.}  Non-unique finite outputs remain finite relations unless uniqueness, a gap, or a shrinking condition justifies selection.
\item \textbf{Quantifier-free execution invariants.}  All finite searches, folds, traces, and dynamic programs are represented by Boolean state and proved by quantifier-free induction.
\item \textbf{Stable elimination.}  If the final theorem is independent of intermediate choices, the compiler should eliminate those choices directly rather than manufacture a global selector.
\item \textbf{Explicit quantitative guards.}  Every discontinuous or partial operation states the margin it consumes.
\item \textbf{Producer/verifier separation.}  Deep information such as a height bound, zero-free region, root separation, or finite cover may be an input certificate.  The generic layer proves that a valid certificate suffices; it does not claim to produce unavailable data.
\item \textbf{A short Lean path.}  The proof architecture should reduce to a few reusable modules and exploit the existing object recursion, contract, and observer combinators.
\end{enumerate}

The PFIG framework introduced below meets these goals by adding one decisive normalization result: a fixed graph is compiled to one query-specialized encoded loop.  This removes the need for a highly dependent, node-by-node object-language graph compiler and makes the formalization path much more direct.

\subsection{What ``resolve'' means in this report}

Three outcomes must be distinguished.

\begin{description}
\item[Resolved.] A reusable theorem gives the desired construction from precisely stated finite data, and its proof uses only principles available in iPRA.
\item[Reformulated.] The original unrestricted statement asks for a discontinuous or noncomputable selector.  A stable relational, projective, clustered, or guarded replacement is proved instead, together with a barrier explaining why the unrestricted version should not be expected.
\item[Isolated.] The framework reduces the problem to a named quantitative producer, such as an explicit height bound.  The finite verifier and all subsequent construction are resolved, but the producer remains theorem-specific or open.
\end{description}

This vocabulary prevents a common form of overstatement: a framework can genuinely solve the logical and representational part of a theorem while leaving its deepest quantitative input untouched.

\section{Logical base and proof-level discipline}
\label{sec:logical-base}

\subsection{Three levels that must not be conflated}

The framework uses three levels of description.  Keeping them separate is the first correctness invariant.

\paragraph{Level 0: metatheory.}
Lean proves facts about syntax, substitutions, lists, finite graphs, and derivation constructors.  At this level we may recurse structurally over a fixed graph and build a derivation tree.  Such recursion is not asserted to occur inside iPRA.

\paragraph{Level 1: object derivability.}
A theorem of the form
\[
  h : \Prf\ \PRAAx\ \Gamma\ \varphi
\]
is a syntactic iPRA derivation.  The constructors in \texttt{Deduction.lean} implement intuitionistic natural deduction with equality.  The derived lemmas for substitution, weakening, cut, and equality congruence allow these derivations to be assembled compositionally \cite{repo-deduction}.

\paragraph{Level 2: external semantics and reference execution.}
A Python program or a semantic Lean structure may evaluate a finite relation, test a dynamic program, or illustrate an error estimate.  This is useful evidence and a specification oracle, but it becomes an iPRA proof only after the corresponding primitive-recursive symbols, object formulas, and derivation constructors have been supplied.

A file that proves a Level-2 statement and merely names the desired Level-1 theorem has not completed the formalization.  Conversely, a Level-1 proof transformer may be perfectly general while accepting theorem-specific arithmetic estimates as inputs.  This is not a defect: it is the intended producer/verifier boundary.

\subsection{The ordinary object language}

We use the checked ordinary iPRA layer rather than programming directly in the smallest native calculus.  Let \(\PTm_n\) be the primitive-recursive terms with \(n\) free variables and let \(\Fm_n\) be first-order formulas over those terms.  For a context \(\Gamma\subseteq\Fm_n\),
\[
  \Gamma\vdash_{\PRAAx}\varphi
\]
abbreviates \(\Prf\ \PRAAx\ \Gamma\ \varphi\).

The proof system has the expected intuitionistic rules for \(\bot,\wedge,\vee,\to,\forall,\exists\), equality reflexivity, simultaneous Leibniz substitution, and universally instantiated theory axioms.  The key structural facts are:

\begin{align*}
\text{Substitution: }&
  \Gamma\vdash\varphi
  \Longrightarrow
  \Gamma[\sigma]\vdash\varphi[\sigma],\\
\text{Weakening: }&
  \Gamma\vdash\varphi,\ \Gamma\subseteq\Delta
  \Longrightarrow
  \Delta\vdash\varphi,\\
\text{Simultaneous cut: }&
  \Gamma\vdash\varphi,\quad
  (\forall\psi\in\Gamma)(\Delta\vdash\psi)
  \Longrightarrow
  \Delta\vdash\varphi.
\end{align*}

These are already proved by induction on derivations.  PFIG composition therefore does not require a new logic.  It needs a disciplined way to turn each numerical interface into formulas and then to feed its local proofs to these structural lemmas.

\subsection{Quantifier-free induction as a programming rule}

For a quantifier-free formula \(I(t,\vec{x})\), iPRA proves
\[
  I(0,\vec{x})\ \wedge\
  \forall t\bigl(I(t,\vec{x})\to I(t+1,\vec{x})\bigr)
  \to \forall t\,I(t,\vec{x}).
\]
At the object level, the repository provides a derived rule \texttt{PRA.qfInd}.  The existing \path{ObjectIteration.lean} theorem goes further: if a state is computed by primitive recursion and \(P(s,t,\vec{x})\) is a quantifier-free state invariant, then object derivations of the base and local update yield an object derivation of the invariant for the computed state at every time, and hence at any supplied fuel term \cite{repo-local-global}.

This theorem should be treated as a verified loop rule.  It says that an ordinary proof of a finite algorithm need only provide:

\begin{enumerate}
\item a fixed primitive-recursive initial-state symbol;
\item a fixed primitive-recursive update symbol;
\item a quantifier-free formula describing the invariant;
\item an object derivation of the initial invariant;
\item an object derivation that one update preserves the invariant.
\end{enumerate}

Everything else---iteration to a variable fuel, specialization to a computed bound, and equality transport between the recursion symbol and the intended state---is generic.

\subsection{Fixed symbols, not a universal interpreter}

A primitive-recursive language may contain one symbol for every primitive-recursive function without containing a single primitive-recursive function that evaluates arbitrary codes for all primitive-recursive functions.  The distinction is essential.

Suppose a mathematical construction uses a particular polynomial evaluator, a particular finite-state transition, and a particular Boolean certificate checker.  The metatheory generates three fixed \(\PRA\) terms for these functions.  Their defining equations are available as theory axioms.  A generic Lean compiler may inspect the finite syntax of the construction and output these terms.  Nothing in iPRA is asked to interpret an arbitrary program code.

Alternatively, if a program family is data, one may use a \emph{clocked trace relation}.  The checker verifies a proposed trace of a fixed small machine for a supplied number of steps.  Trace verification is primitive recursive even when unbounded evaluation is not total.  The PFIG normal form uses both patterns:

\begin{itemize}
\item fixed mathematical libraries are statically specialized to fixed PR symbols;
\item imported computations may be accepted through finite traces whose local transitions are checked.
\end{itemize}

\subsection{Formulas as interfaces}

A PFIG node is not merely a function.  It carries an interface with four logically distinct components.

\begin{definition}[Finitary interface]
For input parameters \(\vec{x}\), a precision/resource request \(p\), and an output code \(y\), a finitary interface consists of:
\begin{enumerate}
\item a quantifier-free validity predicate \(D(\vec{x})\);
\item a primitive-recursive bound \(B(p,\vec{x})\);
\item a quantifier-free answer relation \(R(y,p,\vec{x})\), used only with \(y<B(p,\vec{x})\);
\item a demand transformer \(d(p,\vec{x})\) describing the input precision needed to certify an output at request \(p\).
\end{enumerate}
The local totality contract is
\[
  D(\vec{x})\to \exists y<B(p,\vec{x})\ R(y,p,\vec{x}).
\]
An exact interface additionally supplies a term \(w(p,\vec{x})\) and a derivation of \(R(w(p,\vec{x}),p,\vec{x})\).
\end{definition}

The bound is not cosmetic.  It turns existential search, universal checking, and case analysis into primitive-recursive operations.  The relation is not cosmetic either: it allows honest representation of non-unique outputs.

\subsection{Three elimination modes}

Every internal output is consumed in one of three ways.

\begin{description}
\item[Exact selection.] A primitive-recursive term produces a witness, and its correctness proof is substituted directly into the consumer.
\item[Finite search.] A bounded decidable relation is known nonempty.  The least valid code is found by bounded search, and a reflection theorem proves that it satisfies the relation.
\item[Stable erasure.] The consumer proves the same terminal statement for every valid internal witness.  Existence plus this local stability proves the terminal statement; no selector is constructed.
\end{description}

These modes subsume most apparent differences among the repository packets.  For example, deterministic arithmetic is exact selection; a finite list of candidate roots is finite search; and an eigenvalue multiset identity that is independent of eigenvector choices uses stable erasure.

\subsection{Object contracts}

The simplest exact object contract is already present in the repository.  For precondition \(P(\vec{x})\) and postcondition \(Q(y,\vec{x})\), it consists of a term \(w(\vec{x})\) and a derivation
\[
   P(\vec{x})\vdash_{\PRAAx} Q(w(\vec{x}),\vec{x}).
\]
Substitution composes such contracts.  If the second contract is proved under \(Q\), instantiate its free output variable by \(w\), then cut with the first certificate.  This is exactly the construction implemented in \texttt{ObjectComposition.lean} \cite{repo-coherent}.

The general framework will retain this exact core and add bounded-search and stable-erasure adapters.  Crucially, each adapter is proved once.  A theorem-specific package should not reimplement bounded induction, least search, witness elimination, or precision bookkeeping.

\subsection{Correctness criterion for a completed embedding}

A mathematical theorem has been embedded in iPRA only when all of the following are explicit.

\begin{enumerate}
\item The input and output objects are represented by natural-number codes or primitive-recursive names.
\item Validity, finite observations, and certificate checking are primitive-recursive and represented by object formulas.
\item Every use of induction has a displayed quantifier-free invariant.
\item Every existential output is exact, bounded-searchable, or eliminated by stability.
\item Every partial or discontinuous operation has a quantitative guard or a relational replacement.
\item The final Lean theorem has conclusion \(\Prf\ \PRAAx\ \Gamma\ \varphi\), not merely a semantic proposition about standard naturals.
\item Any transport to the native scalar core is applied only after item 6.
\end{enumerate}

This criterion is deliberately stronger than ``the algorithm is primitive recursive.''  Primitive recursiveness of the data path is necessary, but the proof path must also be represented and composed.

\section{Proof-producing finitary interface graphs}
\label{sec:pfig}

\subsection{Graph syntax}

A proof-producing finitary interface graph, abbreviated PFIG, is a fixed finite directed acyclic graph.  Vertices are interfaces; edges carry numerical codes and precision demands.  ``Fixed'' means that the number and kinds of vertices belong to the theorem statement and are known to the metatheory.  Inputs such as a polynomial, a matrix dimension, a precision, or a finite complex remain variables inside the object theory.

\begin{definition}[PFIG node kinds]
A node is one of the following.

\begin{enumerate}
\item \textbf{Exact transducer.}  A fixed primitive-recursive term \(f\), with an object proof that its output satisfies a postcondition.
\item \textbf{Bounded loop.}  A primitive-recursive initial state, step function, fuel, and quantifier-free invariant, with object proofs of initialization, preservation, and terminal extraction.
\item \textbf{Bounded search.}  A primitive-recursive bound and quantifier-free predicate, together with an object proof of nonemptiness when the precondition holds.
\item \textbf{Relational producer.}  A finite answer cloud \(R(y,p,\vec{x})\) with a bound and a proof that at least one valid answer exists.  The node does not promise a coherent selector across precisions.
\item \textbf{Guarded operation.}  A total primitive-recursive operation whose intended algebraic law is proved under a positive numerical margin, such as \(|x|>2^{-g}\).
\item \textbf{Stable observer.}  A consumer that proves a terminal formula \(O\) from the validity of one or more internal certificates, independently of their exact codes.
\item \textbf{Finite gluing node.}  A bounded tuple of local codes and a quantifier-free overlap checker, with componentwise constructions and a local-to-global extraction proof.
\end{enumerate}
\end{definition}

Every node specifies a backward demand map.  If the node must produce output precision \(p\), its incoming edge \(e\) requests precision \(d_e(p,\vec{x})\).  A node with several consumers receives the maximum of their requests.  Since the graph is finite and acyclic, all demands are computed by a reverse topological pass.

\subsection{Why a graph rather than a single representation}

No one data type is appropriate for all mathematical concepts.  Exact finite algebra is best represented by canonical integer or rational codes.  Continuous objects are best represented by regular names or finite observations.  Roots and eigenspaces are naturally finite relations near collisions.  Local geometric or sheaf data are naturally finite compatible families.  Proofs themselves often mix all of these.

A graph makes this heterogeneity explicit while maintaining one compilation target.  A node may, for example, take an exact polynomial code, call a bounded subdivision loop, return a finite cloud of isolating boxes, pass those boxes to a symmetric coefficient observer, and finally produce a rational inequality.  The output of the whole graph is not required to have the same representational mode as any intermediate node.

\subsection{Requests and active state}

Fix a terminal request \(p\) and input \(\vec{x}\).  The reverse pass computes a demand \(r_v(p,\vec{x})\) at every vertex \(v\).  The forward pass then computes only the finite data needed at those requests.

For a topological ordering \(v_0,\ldots,v_{N-1}\), define the active state after stage \(i\) to contain:

\begin{itemize}
\item the request values \(r_{v_j}\) for all relevant vertices;
\item all exact outputs already computed;
\item finite tables or traces required by later nodes;
\item for relational nodes, either a bounded answer table and viability bits, or only a certificate sufficient for a stable observer;
\item a Boolean flag for each local invariant and compatibility condition.
\end{itemize}

All tuples and finite tables are encoded by standard primitive-recursive pairing and bounded-sequence encodings.  Because the graph is fixed, the tuple shape is fixed.  Decoding an active state is therefore a fixed family of primitive-recursive projections, not a universal data interpreter.

\subsection{The specialization discipline}

The graph compiler is a metatheoretic function on the finite PFIG syntax.  Given a particular graph \(G\), it emits:

\begin{enumerate}
\item a fixed primitive-recursive demand function \(\mathsf{demand}_G\);
\item a fixed primitive-recursive initial state \(\mathsf{init}_G\);
\item a fixed primitive-recursive transition \(\mathsf{step}_G\);
\item a primitive-recursive number of stages \(N_G\), or a request-dependent bounded fuel when a node contains a finite loop;
\item a quantifier-free state invariant \(\mathsf{Inv}_G\);
\item object derivations of the base, step, and final extraction formulas, assembled from local certificates.
\end{enumerate}

The resulting object theorem quantifies over \(p\) and \(\vec{x}\), but not over arbitrary graph syntax or arbitrary program codes.  This is the exact sense in which PFIGs are general without requiring an internal universal evaluator.

\subsection{Graph normal form}

The key simplification is that a fixed PFIG need not be formalized by a dependent object-language interpreter for graph nodes.  It can be normalized to a small number of ordinary primitive recursions.

\begin{theorem}[PFIG normal form]
\label{thm:pfig-normal-form}
Let \(G\) be a fixed finite acyclic graph whose local operations, bounds, demand maps, and certificate checkers are primitive recursive.  Assume that every local correctness obligation is supplied as an object-iPRA derivation, and every relational node is consumed by finite search, finite-horizon viability, or stable erasure.  Then one can effectively construct:

\begin{enumerate}
\item a fixed primitive-recursive encoded state transition \(S_G\);
\item a primitive-recursive fuel \(F_G(p,\vec{x})\);
\item a quantifier-free invariant \(I_G(s,t,p,\vec{x})\);
\item object-iPRA derivations of the base and preservation sequents for \(I_G\);
\item a terminal extraction derivation from \(I_G\) at fuel \(F_G\).
\end{enumerate}

Consequently, the existing proof-producing primitive-recursion theorem yields an object derivation of the terminal PFIG contract.
\end{theorem}

\begin{proof}
Choose a topological ordering of the graph.  We first normalize each node to one or more finite instructions of four elementary kinds:

\begin{enumerate}[label=(\alph*)]
\item assign a fixed primitive-recursive expression to a state field;
\item update a bounded-loop state once;
\item update one cell of a bounded search or viability table;
\item set a Boolean certificate flag after checking a quantifier-free relation.
\end{enumerate}

An exact transducer gives an instruction of kind (a).  A bounded loop contributes its own finite sequence of kind-(b) instructions.  A bounded search is a left-to-right fold that stores whether a witness has been seen and, if so, the least one.  A relational finite-horizon computation contributes a finite table whose cells are filled in reverse level order; each cell uses bounded existential search over the next level.  A stable observer contributes no choice: it contributes a final implication whose premises are certificate flags.  A guarded operation is an exact assignment accompanied by a flag expressing its margin.  A finite gluing node is a bounded loop over charts or overlaps.

The reverse demand pass is itself a fixed reverse topological fold.  Because \(G\) is fixed, its adjacency lists and edge transformers are compiled into the transition syntax.  At each instruction, the current request field is replaced by a maximum of finitely many primitive-recursive expressions.  Maximum, pairing, projections, bounded indexing, and bounded table update are primitive recursive.

Concatenate the reverse-demand instructions and the forward-execution instructions into one schedule.  Encode the schedule position \(t\) and all active fields in one natural number \(s\).  Define \(S_G(s,t,p,\vec{x})\) by a finite case distinction on the schedule position.  Since the number of outer instruction schemas is fixed, this case distinction is compiled to primitive recursion or nested selectors.  Inner loops use counters and bounds stored in the state.  A single global fuel is obtained by summing the primitive-recursive lengths of all scheduled segments.  Thus both \(S_G\) and \(F_G\) are primitive recursive.

The invariant \(I_G\) is the conjunction of the following quantifier-free clauses, encoded as one Boolean checker if desired:

\begin{itemize}
\item the schedule position and segment counters are within their bounds;
\item every completed demand field equals the result of its compiled reverse equation;
\item every completed exact field satisfies its local postcondition;
\item every completed loop field satisfies the supplied loop invariant;
\item every completed search field satisfies the standard least-search invariant;
\item every completed viability cell satisfies its defining finite recurrence;
\item every set certificate flag equals one and hence reflects its intended relation;
\item untouched fields retain their previous values.
\end{itemize}

All quantification implicit in ``every completed cell'' is implemented by a bounded Boolean fold stored in the invariant checker.  Hence the displayed object formula is quantifier-free.

For the base sequent, unfold \(\mathsf{init}_G\), the zero equations of the state constructors, and the empty-fold identities.  Each invariant clause is immediate.  For preservation, perform the finite case split on the schedule position.  In an exact assignment case, use the node's local object certificate and equality substitution.  In a loop case, use its supplied one-step proof.  In a search or viability case, use Boolean reflection and bounded arithmetic.  In a stable-observer case, no state choice is required; only the corresponding local implication is recorded.  The metatheoretic compiler builds these cases by repeated use of the object proof constructors and simultaneous cut.

The object primitive-recursion theorem now proves \(I_G\) at every time and, in particular, at \(F_G(p,\vec{x})\).  By the definition of the final schedule position, all required node clauses and certificate flags are available.  Cut them into the terminal extraction derivation.  The conclusion is the desired terminal contract.
\end{proof}

\subsection{Two-loop refinement}

For implementation, it is often clearer to use two loops rather than one:

\begin{enumerate}
\item \(S_G^-\) runs the reverse demand pass;
\item \(S_G^+\) runs the forward certificate pass using the completed demand vector.
\end{enumerate}

The two-loop form produces smaller invariants and simpler Lean terms.  The one-loop theorem is a conceptual normal form, not a mandate to maximize compression.  Since exact term contracts compose, the output state of the first loop is simply substituted into the second loop's precondition.

\subsection{Normal form for relational outputs}

A relational node can be normalized in three different ways, matching the three elimination modes.

\paragraph{Least finite witness.}
At a fixed precision, compute the least \(y<B\) satisfying \(R(y)\).  The state stores \((\mathsf{found},y)\); the invariant says that no smaller valid witness was skipped and that, when \(\mathsf{found}=1\), \(R(y)\) holds.

\paragraph{Finite-horizon path.}
For levels \(0,\ldots,H\), compute a backward viability table, then follow the least viable successor.  Both passes are bounded.  This gives a coherent path only to the requested finite horizon.

\paragraph{Stable erasure.}
Do not select an internal code.  Prove
\[
  R_1(y_1)\wedge\cdots\wedge R_k(y_k)\to O
\]
for arbitrary valid codes, combine this implication with the independently produced certificates, and discard the codes.  This is often the mathematically strongest and computationally cheapest option.

\subsection{Composition theorem}

\begin{theorem}[PFIG composition]
\label{thm:pfig-composition}
Suppose \(G_1\) proves a finite interface contract from \(P\) to \(Q\), and \(G_2\) proves a contract from \(Q\) to \(R\), with compatible precision demands.  Then the graph obtained by connecting the output of \(G_1\) to the input of \(G_2\) has a PFIG compilation from \(P\) to \(R\).  The composite demand is the primitive-recursive composition of the two demands; the composite object proof is substitution followed by cut.
\end{theorem}

\begin{proof}
For an exact output, this is the existing term-contract composition theorem.  For a bounded-search output, instantiate the second certificate at the least witness and cut with the search reflection theorem.  For a relational output that is stably erased by \(G_2\), combine existence from \(G_1\) with the stability implication of \(G_2\) by existential elimination.  Precision demands compose because a request \(p\) at the terminal node first induces \(d_2(p)\) at the interface and then \(d_1(d_2(p))\) at the original input.  All maps are fixed primitive-recursive terms, so substitution gives the object formulas at the composite request.
\end{proof}

\subsection{Why this normal form improves the formalization problem}

A direct graph compiler in Lean must manage dependent vectors of node types, formulas with changing arities, substitutions along heterogeneous edges, and a proof that every generated program symbol has the right defining equations.  That is possible, but it front-loads a large amount of engineering before a single mathematical packet benefits.

The normal form separates concerns.  A small external elaborator produces ordinary PR syntax and one ordinary invariant.  The Lean proof layer needs only:

\begin{enumerate}
\item reflection lemmas for Boolean folds and bounded search;
\item the existing primitive-recursion invariant theorem;
\item exact contract composition and stable-observer elimination;
\item a correctness proof for the external elaborator's finite syntax translation, or initially a theorem-specific generated Lean file.
\end{enumerate}

Thus a first end-to-end theorem can be completed without first formalizing a universal PFIG AST.  Once several generated examples are stable, the elaborator itself can be internalized in Lean and proved sound by structural recursion.

\section{Bounded Booleanization, search, and finite quantification}
\label{sec:booleanization}

The single most reusable compilation step is to replace a bounded formula by a primitive-recursive Boolean whose value is proved equivalent to that formula.  This simultaneously solves three problems: bounded quantifiers disappear from induction invariants, finite case distinctions become intuitionistically legitimate, and bounded witnesses can be selected by least search.

\subsection{Boolean conventions}

Represent false by \(0\) and true by \(1=S(0)\).  Define primitive-recursive functions
\[
 \mathsf{not}(b),\quad \mathsf{and}(b,c),\quad \mathsf{or}(b,c),
 \quad \mathsf{isZero}(n),\quad \mathsf{eq}(m,n),\quad \mathsf{lt}(m,n),
\]
with outputs in \(\{0,1\}\).  The exact choice of definitions is unimportant provided their equations and range properties are proved in iPRA.  For example, equality can be computed by simultaneous bounded subtraction or by recursion on the two arguments; order can be computed from truncated subtraction.

For a Boolean term \(b\), write \([b]\) for the object formula \(b=1\).  Each Boolean connective has an object-iPRA reflection law, such as
\begin{align*}
 [\mathsf{not}(b)] &\leftrightarrow \neg[b],\\
 [\mathsf{and}(b,c)] &\leftrightarrow [b]\wedge[c],\\
 [\mathsf{or}(b,c)] &\leftrightarrow [b]\vee[c].
\end{align*}
The proof is a finite case distinction on the Boolean range of the arguments.  The range statement itself is quantifier-free and follows by primitive-recursive induction on the defining computation.

\subsection{Bounded formulas}

A bounded formula is generated from atomic equalities and inequalities by \(\wedge,\vee,\to,\neg\), and bounded quantifiers
\[
 \forall i<t\ A(i),\qquad \exists i<t\ A(i),
\]
where all bounds and terms are primitive recursive.  The formula may have arbitrary free numerical parameters, but the syntactic nesting depth is fixed.

\begin{theorem}[Structural Booleanization]
\label{thm:booleanization}
For every fixed bounded formula \(A(\vec{x})\), one can construct a primitive-recursive Boolean term \(\chi_A(\vec{x})\) and an object-iPRA derivation
\[
  \vdash_{\PRAAx} [\chi_A(\vec{x})]\leftrightarrow A(\vec{x}).
\]
Moreover, iPRA proves \(\chi_A(\vec{x})\le 1\).
\end{theorem}

\begin{proof}
Proceed by structural induction on \(A\).

For an equality \(s=t\), take \(\chi_A=\mathsf{eq}(s,t)\).  The defining correctness proof for \(\mathsf{eq}\) gives the equivalence.  Other decidable atomic relations are treated similarly.

For conjunction, disjunction, and negation, apply the corresponding Boolean operation and compose the induction hypotheses with its reflection law.  For implication, set
\[
  \chi_{A\to B}=\mathsf{or}(\mathsf{not}(\chi_A),\chi_B).
\]
Because the induction hypothesis supplies decidability of \(A\), intuitionistic logic proves
\(
eg A\vee B\leftrightarrow(A\to B)\).  Explicitly, the forward direction is immediate.  For the reverse direction, case-split on \(A\vee\neg A\); if \(A\), use the implication to obtain \(B\), and if \(\neg A\), take the left disjunct.  The case split is constructive because it is obtained from the Boolean value \(\chi_A\), not from unrestricted excluded middle.

For bounded universal quantification, define a fold
\begin{align*}
 \mathsf{All}_A(0,\vec{x}) &=1,\\
 \mathsf{All}_A(n+1,\vec{x})
  &=\mathsf{and}(\mathsf{All}_A(n,\vec{x}),\chi_A(n,\vec{x})).
\end{align*}
Use quantifier-free induction on \(n\) to prove
\[
 [\mathsf{All}_A(n,\vec{x})]
 \leftrightarrow \forall i<n\ A(i,\vec{x}).
\]
The right side in this displayed metatheoretic statement need not be the induction matrix.  In the actual object proof, the invariant is a quantifier-free equality relating the fold accumulator to the recursively computed Boolean table, or equivalently a Booleanized bounded universal.  The induction step uses only the recursion equation and the already proved reflection law for \(\chi_A\).  Specialize at the bound term \(t(\vec{x})\).

For bounded existential quantification, define
\begin{align*}
 \mathsf{Any}_A(0,\vec{x}) &=0,\\
 \mathsf{Any}_A(n+1,\vec{x})
  &=\mathsf{or}(\mathsf{Any}_A(n,\vec{x}),\chi_A(n,\vec{x})),
\end{align*}
and prove the analogous reflection theorem.  The output-range property is preserved by every Boolean constructor and fold.
\end{proof}

\subsection{A precise quantifier-free induction pattern}

It is worth spelling out how the universal-fold proof avoids placing a quantifier in the induction matrix.  Define an encoded prefix table
\[
 T_A(n,\vec{x})=\langle \chi_A(0,\vec{x}),\ldots,
                    \chi_A(n-1,\vec{x})\rangle
\]
by primitive recursion.  Define \(\mathsf{allTable}(T,n)\) by a bounded fold.  The induction invariant is the literal equation
\[
  a_n=\mathsf{allTable}(T_A(n,\vec{x}),n),
\]
where \(a_n\) is the accumulator generated by the recursion.  It is quantifier-free.  After the recursion theorem has established this equality, a separate finite-fold reflection lemma proves that \(a_n=1\) is equivalent to the bounded universal formula.  The proof therefore never uses induction on a quantified assertion.

The same pattern applies to finite sums, products, maxima, list membership, matrix identities over bounded indices, path checking, and coverage of a finite net.  One should regard a quantified mathematical statement as the \emph{specification} of a Boolean fold, not as the invariant itself.

\subsection{Least bounded search}

Let \(A(i,\vec{x})\) be bounded-decidable with Boolean \(\chi_A\).  Define a search state \((f,w)\), where \(f\in\{0,1\}\) records whether a witness has been found and \(w\) stores the first witness.  Initialize \((0,0)\).  At stage \(i\):
\[
 (f',w')=
 \begin{cases}
 (1,i),& f=0\text{ and }\chi_A(i,\vec{x})=1,\\
 (f,w),&\text{otherwise.}
 \end{cases}
\]
This is primitive recursive.

\begin{theorem}[Least bounded search]
\label{thm:least-search}
There is a primitive-recursive function \(\mu_{<B}A(\vec{x})\) such that iPRA proves:

\begin{enumerate}
\item if \(\exists i<B\ A(i,\vec{x})\), then
\(\mu_{<B}A<B\) and \(A(\mu_{<B}A,\vec{x})\);
\item if \(j<\mu_{<B}A\), then \(\neg A(j,\vec{x})\);
\item if no witness exists, the function returns a fixed default, say \(0\).
\end{enumerate}
\end{theorem}

\begin{proof}
Run the search state for \(B\) steps.  Use the quantifier-free invariant consisting of the following Boolean equalities:

\begin{itemize}
\item \(f\le 1\);
\item if \(f=1\), then \(w<i\) and \(\chi_A(w)=1\);
\item the Boolean fold \(\mathsf{Any}_A(i)\) equals \(f\);
\item a bounded Boolean fold confirms that every \(j<w\) has \(\chi_A(j)=0\).
\end{itemize}

The base case follows from the empty-fold equations.  The step is a finite Boolean calculation.  At time \(B\), nonemptiness implies \(\mathsf{Any}_A(B)=1\), hence \(f=1\), and the remaining invariant clauses yield validity and minimality.  Define \(\mu_{<B}A\) to be \(w\) when \(f=1\) and \(0\) otherwise.
\end{proof}

\subsection{Finite choice and its exact scope}

Bounded search yields a form of finite choice strong enough for most PFIG nodes.

\begin{corollary}[Bounded finite choice]
Suppose iPRA proves
\[
  \forall x<C\ \exists y<B(x)\ R(x,y)
\]
where \(R\) is bounded-decidable.  Then a primitive-recursive function \(f(x)\) can be defined by least bounded search, and iPRA proves \(R(x,f(x))\) for every \(x<C\).
\end{corollary}

There is no conflict with constructive restrictions: the domain and every search range are explicitly finite and decidable.  What does \emph{not} follow is a coherent selector from an infinite sequence of unrelated finite nonempty sets.  Coherence across precision is a separate theorem requiring compatibility and shrinking hypotheses, treated in \cref{sec:shrinking}.

\subsection{Booleanized certificates}

A certificate checker is a primitive-recursive Boolean \(\mathsf{check}(c,\vec{x})\).  The reflection theorem has the form
\[
 [\mathsf{check}(c,\vec{x})]\to \mathsf{Specification}(c,\vec{x}).
\]
The specification may contain quantifiers; only the checker and loop invariant must be quantifier-free.  This separation is useful for algebraic identities, paths in finite graphs, polynomial factorizations, triangulations, covering lists, determinant expansions, and finite analytic estimates.

A proof-producing checker node therefore needs two object derivations:

\begin{enumerate}
\item the finite computation establishes \(\mathsf{check}=1\);
\item reflection turns \(\mathsf{check}=1\) into the mathematical postcondition.
\end{enumerate}

The second proof is often theorem-specific but structurally simple.  It unfolds the certificate format and uses arithmetic lemmas; it does not repeat the search or trace argument.

\subsection{Compilation of finite sums and products}

For later applications, let \(a(i,\vec{x})\) be a primitive-recursive integer or rational-code expression.  Define
\begin{align*}
 \mathsf{Sum}(0)&=0,&
 \mathsf{Sum}(n+1)&=\mathsf{Sum}(n)+a(n),\\
 \mathsf{Prod}(0)&=1,&
 \mathsf{Prod}(n+1)&=\mathsf{Prod}(n)\cdot a(n).
\end{align*}
Their algebraic laws are proved by quantifier-free induction on equality invariants.  A nested finite sum is just a nested primitive recursion.  Uniform dimension does not create higher-type difficulty: matrices are bounded tables, matrix multiplication is a bounded sum, and determinant certificates can be checked by a bounded permutation sum or a division-free recurrence.

\subsection{Bounded minimization and finite nets}

If \(v(i)\) is a rational code for \(i<N\), compare values by exact cross multiplication and run a least-argmin fold.  iPRA proves that the selected index is within range and no candidate has smaller value.  If a finite net approximates a compact presentation with modulus \(m(p)\), then approximate minimization follows by:

\begin{enumerate}
\item compute the finite \(2^{-m(p)}\)-net;
\item evaluate a rational approximation to the objective at every net point;
\item choose the least approximate minimizer;
\item use the supplied continuity modulus to transfer the finite comparison to the represented object.
\end{enumerate}

The compactness content is entirely in the finite-net producer and its coverage certificate.  The minimization compiler is generic.

\subsection{Interface theorem for bounded logic}

Combining the preceding results gives the reusable statement needed by the graph compiler.

\begin{theorem}[Bounded-logic interface]
\label{thm:bounded-interface}
For every fixed bounded formula family and every fixed finite collection of folds, searches, sums, products, maxima, or table checks built from it, there are fixed primitive-recursive terms and object-iPRA derivations of their reflection, range, and least-witness laws.  All loop invariants used to establish these laws are quantifier-free.
\end{theorem}

This theorem is schematic in the paper and should become a syntax-directed Lean compiler.  Its syntax is intentionally limited to bounded formulas; unbounded quantifiers remain in specifications and are never evaluated by the compiler.

\section{Proof-producing bounded traces}
\label{sec:traces}

A bounded computation becomes useful to iPRA only when its execution proof is internalized.  The repository already contains the essential object theorem for primitive recursion.  This section packages it as a trace-reflection interface and shows how to incorporate imported finite computations without asserting an internal universal evaluator.

\subsection{Deterministic state recursion}

Fix parameters \(\vec{x}\).  Let \(f(\vec{x})\) be an initial state and let
\[
 g(t,s,\vec{x})
\]
be a primitive-recursive update.  Define
\begin{align*}
 \mathsf{run}(0,\vec{x}) &= f(\vec{x}),\\
 \mathsf{run}(t+1,\vec{x}) &= g(t,\mathsf{run}(t,\vec{x}),\vec{x}).
\end{align*}
Let \(P(s,t,\vec{x})\) be quantifier-free.

\begin{theorem}[Object trace invariant]
\label{thm:trace-invariant}
Suppose there are object-iPRA derivations
\begin{align*}
 \Gamma(\vec{x}) &\vdash P(f(\vec{x}),0,\vec{x}),\\
 P(s,t,\vec{x}),\Gamma(\vec{x})
 &\vdash P(g(t,s,\vec{x}),t+1,\vec{x}).
\end{align*}
Then, for every primitive-recursive fuel term \(F(\vec{x})\), there is an object-iPRA derivation
\[
 \Gamma(\vec{x})\vdash
 P(\mathsf{run}(F(\vec{x}),\vec{x}),F(\vec{x}),\vec{x}).
\]
\end{theorem}

\begin{proof}
This is the mathematical content of the existing \texttt{primitiveRecursionInvariant} and \texttt{primitiveRecursionAt} theorems.  Apply quantifier-free induction to the formula
\[
 P(\mathsf{run}(t,\vec{x}),t,\vec{x}).
\]
It is quantifier-free because \(P\) is, and substitution of PR terms preserves quantifier-freeness.  The base case transports the supplied base proof along the zero equation for primitive recursion.  For the step, substitute \(s=\mathsf{run}(t,\vec{x})\) into the local preservation proof, then transport along the successor equation for primitive recursion.  Universal elimination at \(F(\vec{x})\) gives the conclusion.  The Lean implementation performs the necessary de Bruijn substitutions and equality congruence explicitly \cite{repo-local-global}.
\end{proof}

\subsection{Terminal extraction}

A loop normally proves a state invariant, while the desired theorem is a property \(C(\vec{x})\).  Supply a final derivation
\[
  \Gamma\vdash P(\mathsf{run}(F),F,\vec{x})\to C(\vec{x}).
\]
Modus ponens with 
\cref{thm:trace-invariant} yields \(C\).  The repository's stable-observer module calls this pattern \texttt{loopProducesCertificate} \cite{repo-observable}.

This design avoids placing a complex specification into the induction matrix.  The loop invariant records only finite state facts.  Quantifiers, limits, or abstract mathematical conclusions appear in the terminal extraction theorem, where ordinary intuitionistic reasoning is available.

\subsection{Multiple loop phases}

Many constructions naturally have several phases: build a table, scan it, refine candidates, and verify a certificate.  There are two equivalent encodings.

\paragraph{One tagged state.}
Add a phase field to the state and define a transition by a fixed case distinction.  The invariant is a conjunction of phase-conditional clauses.

\paragraph{Contract composition.}
Treat each phase as a loop contract.  The terminal state of one loop is an exact term in the precondition of the next.  Substitute and cut.  This yields smaller local invariants and is usually preferable in Lean.

The PFIG normal form permits either.  The reference compiler uses explicit phases because they are convenient to inspect, while the Lean roadmap recommends separate loop modules followed by contract composition.

\subsection{Finite trace certificates}

Sometimes the producer is not represented by a simple PR recurrence in the object theory.  It may be an external exact-arithmetic program, a SAT solver, a polynomial factorization routine, or a generated symbolic calculation.  The safe interface is a finite trace certificate.

Let \(\mathsf{Local}(c_i,c_{i+1},i,\vec{x})\) be a quantifier-free relation checking one transition, and let \(T\) encode the list \((c_0,\ldots,c_F)\).  Define a primitive-recursive checker
\[
 \mathsf{TraceOK}(T,F,\vec{x})=
 \mathsf{InitOK}(c_0,\vec{x})\wedge
 \bigwedge_{i<F}\mathsf{Local}(c_i,c_{i+1},i,\vec{x})
 \wedge \mathsf{FinalOK}(c_F,\vec{x}).
\]
Every operation here is a bounded fold over a coded list.

\begin{theorem}[Trace reflection]
\label{thm:trace-reflection}
Suppose object-iPRA proves reflection lemmas for \(\mathsf{InitOK}\), \(\mathsf{Local}\), and \(\mathsf{FinalOK}\).  Then object-iPRA proves
\[
 [\mathsf{TraceOK}(T,F,\vec{x})]	o C(T,F,\vec{x}),
\]
where \(C\) is the terminal specification obtained by composing the reflected local transitions.
\end{theorem}

\begin{proof}
Decode the current trace position and state by bounded projection.  Run a loop over \(i\le F\) whose invariant says that the decoded prefix is locally valid and that the semantic state relation holds at \(c_i\).  The checker equality \(\mathsf{TraceOK}=1\) supplies each local Boolean fact by bounded-table reflection.  Apply the local transition derivation in the step.  At \(F\), use \(\mathsf{FinalOK}\) to derive \(C\).  The invariant is quantifier-free after all prefix checks are represented by Boolean folds.
\end{proof}

This theorem does not assert that the external producer is correct.  It asserts that any accepted trace certifies the result.  Trust is reduced to the small primitive-recursive checker and its object reflection proof.

\subsection{Proof-producing transducers}

A more modular node returns both output data \(y\) and a certificate \(c\).  Its finite checker validates
\[
 \mathsf{check}(c,y,p,\vec{x})=1.
\]
Reflection derives the postcondition \(R(y,p,\vec{x})\).  Such a node may itself be deterministic or relational:

\begin{itemize}
\item deterministic: a fixed PR program computes \((y,c)\);
\item externally certified: the graph receives \((y,c)\) as an input and verifies it;
\item relational: a bounded search enumerates pairs \((y,c)\) whose checker accepts.
\end{itemize}

The downstream graph depends only on \(R\), not on how the certificate was produced.

\subsection{The base, step, final triad}

Every bounded algorithm should document exactly three object-level obligations.

\begin{description}
\item[Base.] The initial encoded state satisfies the invariant.
\item[Step.] Under the invariant, one encoded transition preserves it.
\item[Final.] The terminal invariant implies the advertised contract.
\end{description}

A recurring weakness in draft formalizations is that the base and step are proved semantically in Lean while the final target is merely declared as \(\Prf\ \PRAAx\).  The PFIG discipline requires all three to have the object type.  Semantic proofs may be used to discover the arithmetic argument, but the artifact ledger records the object theorem separately.

\subsection{Trace composition}

Suppose trace \(T_1\) establishes an intermediate relation \(Q(y,\vec{x})\), and trace \(T_2\), parameterized by \(y\), establishes \(R(z,y,\vec{x})\).  The combined trace concatenates the two coded lists with a phase tag.  Reflection can be proved either by a new loop or, more simply, by:

\begin{enumerate}
\item reflect \(T_1\) to obtain \(Q(y,\vec{x})\);
\item substitute \(y\) into the reflection theorem for \(T_2\);
\item cut the first proof into the second.
\end{enumerate}

This is why exact object contracts are the preferred boundary between phases.  Concatenating raw traces is an optimization, not a logical necessity.

\subsection{Nondeterministic traces}

A finite trace relation may admit several next states.  Do not immediately choose the least successor.  At a fixed horizon, one may compute viability backward and then choose a path, as in \cref{sec:viability}.  If the final property is stable under all accepted traces, use stable erasure and avoid path selection entirely.

For example, a polynomial factorization checker may accept several orderings of irreducible factors.  A theorem about their product is stable under order and needs no canonical factor list.  A theorem referring to the ``first factor'' requires either a canonical sorting rule with decidable comparison or an explicit statement invariant under permutation.

\subsection{Resource bounds}

The fuel \(F(p,\vec{x})\), trace length, table sizes, and intermediate bit lengths are part of the interface.  iPRA does not require feasible complexity, but explicit bounds serve three purposes:

\begin{enumerate}
\item they make every search and fold primitive recursive;
\item they make validity a bounded formula and hence Booleanizable;
\item they prevent hidden uses of unbounded termination arguments.
\end{enumerate}

A termination proof by a lexicographic or weighted rank is compiled to a primitive-recursive fuel by bounding the rank range and proving that every nonterminal step decreases it.  The loop may continue with an identity transition after termination; totalization removes partiality from the state function.

\subsection{Trace compiler interface}

The generic Lean module should expose a structure morally equivalent to:

\begin{lstlisting}[language={},caption={Source-level shape of a proof-producing loop contract}]
structure LoopContract (Gamma : List (Fm PTm n)) where
  init      : PR n
  step      : PR (n + 2)
  fuel      : PTm n
  invariant : Fm PTm (n + 2)  -- state, time, parameters
  qf        : invariant.QF
  base      : Gamma |- invariant(init(params), 0, params)
  preserve  : invariant(state,time,params), Gamma
              |- invariant(step(time,state,params), time+1, params)
  result    : Gamma |- invariant(run(fuel),fuel,params) -> post
\end{lstlisting}

The actual Lean syntax uses de Bruijn substitutions and the existing \(\PTm\), \(\Fm\), \(\Prf\), and \(\PRAAx\) types.  The important point is that \texttt{base}, \texttt{preserve}, and \texttt{result} are derivations, not semantic propositions.

\subsection{Resolution achieved}

The trace layer resolves the repository's most pervasive formal gap: how to turn a finite exact computation into an actual iPRA derivation.  It does not prove theorem-specific arithmetic estimates.  It gives a reusable route from those local estimates to a global proof, and it makes the remaining obligations visible and auditable.

\section{Backward demand and the graph compiler}
\label{sec:demand}

Continuous and approximate constructions fail in weak arithmetic when they are formulated ``at all precisions at once.''  The PFIG compiler instead starts with one terminal request, propagates a finite demand backward, and only then executes the finite forward construction.  This section proves the demand calculus and the soundness of the two-pass compiler.

\subsection{Precision contracts}

Let an input presentation have an observation operation \(\Obs_X(x,q)\), interpreted as finite information about \(x\) at precision \(q\).  A node \(F:X\to Y\) has a precision contract if it supplies a primitive-recursive demand map \(d_F(p,\vec{x})\) and proves:
\[
 \Obs_X(x,d_F(p,\vec{x}))\approx
 \Obs_X(x',d_F(p,\vec{x}))
 \quad\Longrightarrow\quad
 \Obs_Y(Fx,p)\approx \Obs_Y(Fx',p).
\]
The equivalence relations may be exact equality, metric closeness, or a finite relational back-and-forth.  The contract is itself represented by a bounded formula or by a stable observer theorem.

For a guarded operation, the demand may also depend on a guard.  For example, to compute \(1/x\) within \(2^{-p}\), one may require an apartness certificate \(|x|\ge 2^{-g}\) and input precision roughly \(p+2g+c\).  The exact bound is theorem-specific, but the composition law is generic.

\subsection{Reverse demand equations}

Let \(G=(V,E)\) be a finite acyclic graph, with a terminal request \(p\) at selected output nodes.  For an edge \(e:u\to v\), let \(d_e(r_v,\vec{x})\) be the request at \(u\) needed by \(v\).  Define requests in reverse topological order:
\[
 r_u=
 \max\Bigl(r_u^{\mathrm{terminal}},
       \max_{e:u\to v} d_e(r_v,\vec{x})\Bigr),
\]
where absent terminal requests are zero.

\begin{theorem}[Primitive-recursive backward demand]
\label{thm:backward-demand}
For every fixed finite acyclic PFIG, the vector \(\vec r(p,\vec{x})\) solving the reverse demand equations is primitive recursive.  iPRA proves each defining equation and the following minimality property: if \(\vec s\) satisfies every terminal and edge demand, then \(r_v\le s_v\) for every vertex \(v\).
\end{theorem}

\begin{proof}
Choose a reverse topological order \(v_{N-1},\ldots,v_0\).  Initialize all request fields with their terminal requests.  At stage \(i\), compute \(r_{v_i}\) as the maximum of its current value and the finitely many expressions \(d_e(r_v,\vec{x})\) for outgoing edges.  Since the graph and each adjacency list are fixed, this is a fixed primitive-recursive fold.  The equality invariant saying that all already processed fields satisfy their defining equation is quantifier-free.  The object proof uses the elementary algebra of maximum.

For minimality, induct in the same reverse order.  Assume all successor requests satisfy \(r_v\le s_v\).  Monotonicity of each demand map gives
\(d_e(r_v)\le d_e(s_v)\le s_u\), and the terminal request is also at most \(s_u\).  Therefore their maximum \(r_u\) is at most \(s_u\).  If a demand map is not monotone as originally written, replace it by its bounded monotone closure
\[
 d_e^*(p)=\max_{q\le p}d_e(q),
\]
which is primitive recursive and satisfies the same correctness contract.
\end{proof}

Minimality is useful but not required for soundness.  It shows that the compiler does not introduce arbitrary excess precision beyond the chosen contract maps.

\subsection{Fix precision, then quantify}

The reverse pass enforces the correct quantifier order.  Suppose a theorem ultimately says
\[
 \forall p\ \exists y\ R(p,y).
\]
The PFIG compiler proves a uniformly primitive-recursive witness or finite relation at each fixed \(p\).  The object proof first treats \(p\) as a free numerical parameter, computes all finite bounds depending on \(p\), runs quantifier-free loops, and only then universally generalizes over \(p\).

It does \emph{not} place ``for every future precision'' in an induction invariant.  This distinction resolves many apparent strength problems.  The algorithm is uniform in \(p\), but every individual execution uses only a finite amount of information.

\subsection{Forward execution}

After the demand vector is known, execute vertices in topological order.  At a vertex \(v\):

\begin{enumerate}
\item read the finite observations requested on incoming edges;
\item run its exact, loop, search, relation, guard, glue, or observer program;
\item store its output code, finite table, or certificate flag;
\item use the local object certificate to extend the global invariant.
\end{enumerate}

No vertex can read an uncomputed successor because of the topological order.  If several consumers need the same output at different precisions, the node is run at their maximum demand; restriction maps or stability prove that this output suffices for each smaller request.

\subsection{Restriction coherence}

A presentation often has restriction maps
\[
 \rho_{q,p}:A_q\to A_p\qquad(q\ge p),
\]
which forget excess precision.  Exact coherence would say
\(\rho_{r,p}=\rho_{q,p}\circ\rho_{r,q}\).  In a relational presentation, coherence may instead be a compatibility relation between restricted clouds.

The graph compiler needs one of the following:

\begin{enumerate}
\item an exact restriction term and object proof of its equations;
\item a stable observer showing that either restriction path gives the same terminal result;
\item a metric estimate bounding the discrepancy and a tail modulus that absorbs it.
\end{enumerate}

Demand sharing is sound under any of these forms.  Requiring literal code equality is unnecessarily strong and often false because codes may contain different but equivalent representatives.

\subsection{Demand for finite relations}

For a relational node, the consumer may ask for more than a numerical precision.  It may need:

\begin{itemize}
\item a horizon \(H\) for compatibility;
\item a diameter target for the viable cloud;
\item a separation margin for exact labels;
\item a finite set of queries rather than one scalar precision.
\end{itemize}

Encode such requests as natural tuples.  The demand order is componentwise or given by an explicit primitive-recursive preorder.  Replace maximum by a primitive-recursive join operation.  The backward theorem then holds verbatim, provided demand maps are monotone with respect to that preorder.

This generalization is important for multiscale arguments.  A consumer of a root cluster may request coefficient precision \(p\), a disk radius bound \(q\), and a multiplicity cutoff \(m\).  Scalarizing everything too early obscures dependencies and produces poor invariants; tuple requests keep them explicit while remaining first-order numerical data.

\subsection{Graph soundness theorem}

\begin{theorem}[Two-pass PFIG soundness]
\label{thm:pfig-soundness}
Let \(G\) be a fixed acyclic PFIG.  Assume:

\begin{enumerate}
\item every request type is naturally coded with primitive-recursive join and order;
\item every edge demand map is primitive recursive and monotone;
\item every node has a local object-iPRA correctness derivation at arbitrary free request parameters;
\item each relational output is consumed by bounded search, finite-horizon viability, or stable erasure;
\item each guarded operation receives the quantitative guard required by its local theorem;
\item the terminal extraction theorem is an object-iPRA derivation from the terminal node contracts.
\end{enumerate}

Then a structural compiler constructs an object-iPRA derivation of the terminal theorem, uniformly in the input and terminal request.
\end{theorem}

\begin{proof}
Compile the reverse pass using \cref{thm:backward-demand} and the trace theorem.  Its final state contains a demand vector satisfying every edge equation.  Substitute these demand terms into all local node contracts.

Compile the forward pass by topological recursion over the finite graph.  Maintain a context containing the postconditions of completed nodes, or equivalently a single conjunction encoded in the state invariant.  For an exact node, apply its local certificate.  For a loop node, use \cref{thm:trace-invariant}.  For a search node, use \cref{thm:least-search}.  For a relational node, invoke the finite viability or stable-erasure theorem.  For a guarded node, cut in the guard certificate.  For a gluing node, use the finite gluing theorem of \cref{sec:gluing}.  At each stage, weakening retains previous contracts and simultaneous cut discharges local assumptions.

At the terminal stage, the compiled context contains every premise of the extraction theorem.  Simultaneous cut produces the terminal derivation.  Universal introduction over the terminal request and input variables gives the advertised uniform theorem.
\end{proof}

\subsection{Branching and conditional nodes}

A finite algorithm may branch on a decidable Boolean.  The graph can represent this as two subgraphs with a merge.  Boolean reflection gives \(B\vee\neg B\), and intuitionistic disjunction elimination combines the branch proofs.  If the branch predicate is an approximate sign without a separation guard, it is not decidable in the represented object and must not be used this way.  The alternatives are:

\begin{itemize}
\item refine until a supplied gap makes the sign decidable;
\item carry both branch outputs as a finite relation;
\item prove the merged observable independent of the branch.
\end{itemize}

Thus the graph syntax itself does not legitimize a discontinuous case split.  Every branch has a Boolean checker whose reflection theorem is part of the node contract.

\subsection{Cycles and iterative refinement}

A mathematical dependency graph may be cyclic, for example when a refinement step updates several mutually dependent approximations.  The PFIG dependency graph remains acyclic by packaging the cycle as one bounded-loop node.  Its state contains all mutually updated components.  This is not a loss of generality: any finitely scheduled iterative process has a primitive-recursive product state.

Unbounded fixed-point iteration is not admitted merely because it converges semantically.  A node must supply a primitive-recursive fuel or a rank bound.  For approximate fixed points, the terminal theorem usually asks only for residual \(\le 2^{-p}\), and the fuel is derived from a contraction or descent modulus.

\subsection{Contract DAG normalization}

A graph with shared subcomputations can be duplicated into a tree, but that obscures resource bounds.  The encoded-state normal form preserves sharing by storing each completed output once.  Demand joining ensures it is computed at sufficient precision for all consumers.  The global invariant records one postcondition per stored node, so proof sharing is reflected as a single local derivation followed by several weakenings.

\subsection{Generated proof terms versus a verified compiler}

There are two practical deployment stages.

\begin{description}
\item[Generated theorem files.]  An external script emits Lean definitions and a proof term specialized to a particular graph.  Lean checks the resulting object derivation.  The generator need not initially be trusted for soundness because malformed output will fail to typecheck; nevertheless, completeness of generation is tested separately.
\item[Internal verified elaborator.]  Define a PFIG syntax in Lean, elaborate it to PR terms and formulas, and prove by structural recursion that the generated derivation has the promised type.  This removes repetition and makes graph-level automation trustworthy.
\end{description}

The first stage is enough to complete repository theorem packets.  The second is a natural consolidation after the interfaces stabilize.

\subsection{What the demand compiler resolves}

The two-pass discipline resolves several reported obstacles at once:

\begin{itemize}
\item all-precision assertions are moved outside induction;
\item only a finite query table is computed for each request;
\item moduli and guards become explicit edge data;
\item shared subcomputations receive a joined demand;
\item exact code equality is replaced by restriction coherence or stable observation;
\item the entire construction remains a fixed family of primitive-recursive symbols.
\end{itemize}

It does not invent missing moduli.  If a theorem supplies only qualitative continuity or ineffective convergence, the corresponding demand edge remains an open producer obligation.

\section{Finite relational presentations and backward viability}
\label{sec:viability}

Many mathematical constructions return a finite set of admissible answers rather than a unique code.  This is not a weakness.  It is often the correct constructive representation.  The problem is to compose such relations without silently assuming countable choice or a continuous labeling that does not exist.

\subsection{Finite answer clouds}

At level or precision \(i\), let
\[
  A_i(\vec{x})=\{a<B_i(\vec{x}) : \Good_i(a,\vec{x})\}
\]
be a finite nonempty set of valid codes.  Between consecutive levels, let
\[
  \Edge_i(a,b,\vec{x})
\]
mean that \(a\in A_i\) and \(b\in A_{i+1}\) are compatible.  All bounds, goodness predicates, and edge predicates are primitive recursive and Booleanized.

Examples include:

\begin{itemize}
\item isolating boxes for roots at increasing precision;
\item approximate eigenspaces or clusters;
\item chart representatives and overlap witnesses;
\item polygonal approximations to a curve;
\item local algebraic branches;
\item finite proof certificates with several valid normal forms.
\end{itemize}

The relation need not have a unique successor or predecessor.  A code may be locally valid but lead to a dead end at a later level.  Therefore choosing the least element of each \(A_i\) independently is unsound.

\subsection{Finite-horizon viability}

Fix a horizon \(H\).  Define viability backward:
\begin{align*}
 \Viable_H^H(a)&:\Longleftrightarrow \Good_H(a),\\
 \Viable_i^H(a)&:\Longleftrightarrow
   \Good_i(a)\ \wedge\
   \exists b<B_{i+1}\,
     \bigl(\Good_{i+1}(b)\wedge\Edge_i(a,b)
                    \wedge\Viable_{i+1}^H(b)\bigr)
\end{align*}
for \(i<H\).  Since every existential is bounded and the recursion has finite depth, the viability table is primitive recursive.

\begin{theorem}[Viability-path equivalence]
\label{thm:viability-path}
For \(i\le H\), \(\Viable_i^H(a)\) holds if and only if there exists a compatible finite path
\[
 a_i=a,a_{i+1},\ldots,a_H
\]
with \(\Good_j(a_j)\) for every \(j\) and \(\Edge_j(a_j,a_{j+1})\) for every \(j<H\).
\end{theorem}

\begin{proof}
Use backward induction on \(H-i\).  At \(i=H\), a path consists only of \(a\), so the statement is exactly the definition of \(\Viable_H^H\).  Suppose the equivalence holds at \(i+1\).  If \(\Viable_i^H(a)\), choose the bounded witness \(b\) in its defining formula.  By the induction hypothesis, \(b\) extends to a path from \(i+1\) to \(H\); prepend \(a\).  Conversely, the second point of any path from \(a\) supplies a witness \(b\), and the remaining suffix implies \(\Viable_{i+1}^H(b)\) by the induction hypothesis.

For the object-iPRA proof, do not induct on the existential formula.  Compute the viability bits in a reverse table.  The invariant states, by Boolean equality, that all completed cells equal the defining bounded Boolean recurrence.  Reflection of the table yields the displayed logical equivalence after the computation is complete.
\end{proof}

\subsection{Least viable path}

If some code at level \(0\) is viable, bounded search can choose the least viable one.  Given a viable code at level \(i\), choose the least successor \(b\) satisfying the edge relation and \(\Viable_{i+1}^H(b)\).  Repeat to the horizon.

\begin{theorem}[Finite dependent choice by viability]
\label{thm:finite-choice}
There is a primitive-recursive function that, from a coded finite relational presentation and horizon \(H\), returns either:

\begin{enumerate}
\item a least lexicographic compatible path of length \(H+1\), or
\item a failure code certifying that no level-zero candidate is viable.
\end{enumerate}

If the input presentation proves nonempty viability, iPRA proves that the returned path is valid and compatible.
\end{theorem}

\begin{proof}
First compute the full viability table by bounded reverse recursion.  Search the level-zero row for the least true cell.  Then move forward: at level \(i\), search the next row for the least \(b\) such that \(\Edge_i(a_i,b)\) and the viability bit of \(b\) is true.  The defining recurrence of the current viability bit proves such a \(b\) exists.  The loop invariant records that the constructed prefix is compatible and that its last point is viable to \(H\).  This invariant is Booleanized and hence quantifier-free.
\end{proof}

This is a genuine finite choice theorem and is fully acceptable in iPRA.  It makes no assertion about a branch through all infinitely many levels.

\subsection{Relational composition}

Suppose \(R(x,y)\) is a finite producer relation and \(S(y,z)\) is a finite consumer relation.  Their relational composite is
\[
 (S\circ R)(x,z)\;:\Longleftrightarrow\;
 \exists y<B_R(x)\bigl(R(x,y)\wedge S(y,z)\bigr).
\]
The existential is bounded, so the composite is again a finite interface.  Nonemptiness follows by existential elimination from nonemptiness of \(R\) and the local totality of \(S\) under each valid \(y\).

For a chain of relations, viability is the efficient implementation of this repeated bounded existential.  It removes intermediate values that cannot reach the final contract.  The graph compiler may either retain the relation as a Boolean table or select a least viable path at the final finite horizon.

\subsection{Backward viability as proof search}

The same dynamic program can be interpreted proof-theoretically.  A candidate is viable when there is a sequence of local certificates leading to a terminal certificate.  The reverse table is therefore a finite cut-elimination aid: it computes which local witnesses can participate in a completed derivation.

This viewpoint is valuable for generated proofs.  The external reference program can output the least viable path, while Lean reconstructs the corresponding object derivation from local proof constructors.  Alternatively, Lean can reflect the whole viability table and use existential elimination.  The latter avoids trusting path extraction but may yield a larger proof term.

\subsection{Dead ends and local validity}

A candidate may satisfy \(\Good_i\) yet be incompatible with every candidate at the next level.  This happens naturally when finite approximations are too coarse or when local certificates omit future constraints.  The framework does not require every good candidate to extend.  It requires either:

\begin{enumerate}
\item at least one viable candidate at each requested finite horizon; or
\item a refined goodness predicate strong enough to imply extension.
\end{enumerate}

The first condition is more flexible.  It allows a producer to overgenerate a finite cloud and lets the consumer's constraints prune it by backward viability.

\subsection{Changing parameters}

Sometimes the relation at level \(i+1\) is computed using parameters refined from level \(i\).  Encode the current parameter code into the state and let
\[
 \Edge_i(a,b,\theta_i,\theta_{i+1})
\]
check both the parameter update and candidate compatibility.  The viability recurrence remains bounded.  A common mistake is to state compatibility while silently holding parameters fixed; the encoded-state formulation makes all changing data explicit.

\subsection{Amalgamation of several witness families}

Suppose a theorem requires simultaneous choices from clouds \(A_i^{(1)},\ldots,A_i^{(k)}\) with cross-compatibility constraints.  Form the product cloud of tuples
\[
 A_i=A_i^{(1)}\times\cdots\times A_i^{(k)}
\]
and include all cross-constraints in \(\Good_i\) and \(\Edge_i\).  Because \(k\) is fixed and each component is bounded, the product bound is primitive recursive.  Backward viability then produces a compatible finite amalgamation.

This handles, for example, simultaneous local branches on a finite cover, a polynomial root together with a multiplicity certificate, or an approximate minimizer together with a subgradient witness.

\subsection{Stable terminal relations}

A consumer may care only that the terminal code lies in a finite equivalence class.  Let \(\sim_H\) be a bounded equivalence or closeness relation on terminal candidates, and suppose the final observable \(O(a_H)\) is invariant under \(\sim_H\).  Then the path selector may be changed without changing the theorem.  The proof should state this invariance and use stable erasure whenever possible.

\subsection{No infinite branch theorem is hidden here}

The existence of a compatible path to every finite horizon does not, constructively and uniformly, provide a single infinite path.  Even classically, extracting such a path invokes a compactness principle; computationally, a family of finite paths may encode noncomputable information.  PFIGs avoid this issue in two ways:

\begin{enumerate}
\item a theorem at fixed precision asks only for a finite path to a primitive-recursive horizon;
\item a coherent regular name is extracted only when shrinking-diameter and motion-tail estimates make the observable value independent of the exact finite path.
\end{enumerate}

The second method is developed in \cref{sec:shrinking}.  It constructs numerical approximants directly; it does not claim to choose an infinite path in the underlying witness tree.

\subsection{Object-language implementation plan}

A Lean module for viability should separate two layers.

\paragraph{Semantic finite layer.}
Define levels as finite lists or bounds with decidable predicates.  Prove the viability-path equivalence, least-path correctness, and product amalgamation in ordinary Lean.  These theorems validate the algorithm and guide testing.

\paragraph{Object reflection layer.}
Represent \(\Good\), \(\Edge\), bounds, table indexing, and table update by fixed PR symbols.  Define a quantifier-free formula saying that a table row satisfies the viability recurrence.  Use the object trace theorem to prove the final table correct.  Reflect true bits to bounded existential formulas.  The final theorem must conclude \(\Prf\ \PRAAx\), not merely the semantic Lean proposition.

The two layers may share definitions of table layout, but their status should be recorded separately in the theorem ledger.

\section{Shrinking viable clouds, coherent names, and stable observers}
\label{sec:shrinking}

Finite-horizon viability solves composition at a fixed precision.  To represent a real number, a root, a compact set, or another infinitary object, one needs a coherent sequence of finite observations.  This section gives sufficient quantitative conditions for constructing such a sequence without selecting an infinite witness branch.

\subsection{Observed values of candidates}

Let \((M,d)\) be a coded metric target, typically rational vectors with an integer-valued scaled distance.  At level \(p\), every candidate \(a\in A_p\) has an observed value
\[
 q_p(a)\in M.
\]
All values and distance comparisons used below are primitive recursive on codes.  Exact real-valued inequalities are represented by rational upper bounds.

Choose a nondecreasing primitive-recursive horizon \(H(p)\ge p\).  Let
\[
 V_p=\{a\in A_p:\Viable_p^{H(p)}(a)\}
\]
be the cloud of candidates that extend to the requested horizon.  Assume \(V_p\) is nonempty; this is a finite certificate obligation.

Define \(a_p\) to be the least code in \(V_p\), and set
\[
 x_p=q_p(a_p).
\]
The choices \(a_p\) need not lie on one path.  Coherence will come from quantitative indistinguishability of viable candidates, not from branch choice.

\subsection{Diameter and edge-motion hypotheses}

Assume primitive-recursive rational bounds \(\del_p\) and \(\eps_p\) with the following properties.

\begin{description}
\item[Viable-cloud diameter.] For all \(a,b\in V_p\),
\[
 d(q_p(a),q_p(b))\le \del_p.
\]
\item[Compatible edge motion.] If \(\Edge_p(a,b)\), then
\[
 d(q_p(a),q_{p+1}(b))\le \eps_p.
\]
\end{description}

The first condition says that all finite witnesses which can survive far enough already agree on the observable value.  The second bounds the change in observation along a compatible refinement.

\begin{theorem}[Adjacent coherence from shrinking viability]
\label{thm:adjacent-coherence}
If \(H\) is nondecreasing, then
\[
 d(x_p,x_{p+1})\le \del_p+\eps_p.
\]
\end{theorem}

\begin{proof}
Because \(a_{p+1}\) is viable to \(H(p+1)\), choose a compatible path from level \(p+1\) to that horizon.  By the definition of viability at level \(p+1\), there is at least one predecessor \(b_p\in A_p\) on a path compatible with \(a_{p+1}\), or, if edges are oriented forward only, take the level-\(p\) member of a viable path whose next point is \(a_{p+1}\).  In implementations where not every viable next-level code has a predecessor, define \(a_{p+1}\) as the least endpoint of a viable pair; the same argument applies.

Since \(H(p+1)\ge H(p)\), truncating this path shows that \(b_p\in V_p\).  Therefore the diameter hypothesis gives
\[
 d(q_p(a_p),q_p(b_p))\le\del_p.
\]
Compatibility of \(b_p\) and \(a_{p+1}\) gives
\[
 d(q_p(b_p),q_{p+1}(a_{p+1}))\le\eps_p.
\]
The triangle inequality yields the claim.
\end{proof}

The minor predecessor convention in the proof is worth making explicit in a formal development.  A clean representation stores viable \emph{pairs} when moving between horizons, so the chosen level-\(p+1\) endpoint always comes with a certified predecessor.

\subsection{Regular-name extraction}

A regular rational name for a real number is often required to satisfy
\[
 |r_p-r_q|\le 2^{-p}\qquad(q\ge p),
\]
or an equivalent adjacent-step condition with a summable tail.

\begin{theorem}[Regular name from a motion budget]
\label{thm:regular-name}
Suppose a primitive-recursive tail modulus proves
\[
  \sum_{j=p}^{q-1}(\del_j+\eps_j)\le 2^{-p}
  \qquad\text{for all }q\ge p.
\]
Then \((x_p)_p\) is a regular Cauchy name.  The function \(p\mapsto x_p\) is primitive recursive.
\end{theorem}

\begin{proof}
Primitive recursiveness follows from bounded viability, least search, and the observation map.  For \(q>p\), repeatedly apply \cref{thm:adjacent-coherence} and the triangle inequality:
\[
 d(x_p,x_q)
 \le\sum_{j=p}^{q-1}d(x_j,x_{j+1})
 \le\sum_{j=p}^{q-1}(\del_j+\eps_j)
 \le 2^{-p}.
\]
Every sum is finite and its inequality is a bounded arithmetic certificate.  The object-iPRA proof uses a finite-sum loop and quantifier-free induction on the equality/inequality code.
\end{proof}

A common convenient schedule is \(\del_p\le 2^{-p-3}\) and \(\eps_p\le 2^{-p-3}\).  Then the tail is at most \(2^{-p-1}\), leaving room for approximation errors in later operations.

\subsection{No infinite witness branch is selected}

The name \(x_p\) is computed afresh from the finite table for horizon \(H(p)\).  The underlying witnesses \(a_p\) may jump between branches.  The diameter hypothesis proves that the observable value is insensitive to this jumping.  Hence the theorem constructs an infinite numerical name without constructing an infinite path in the witness tree.

This distinction resolves a major constructive difficulty.  Root boxes may exchange labels near a collision, local solution branches may bifurcate, and triangulations may refine in noncanonical ways.  If the target observable is stable and the viable clouds shrink, a coherent name still exists.

\subsection{Stable observers}

Often even a regular name is unnecessary.  Suppose two independent producers return certificates \(C_1(c_1)\) and \(C_2(c_2)\), while the desired theorem is \(O\).  A local stability proof has the form
\[
 C_1(c_1),C_2(c_2),\Gamma\vdash_{\PRAAx}O.
\]
Existence or exact production of the certificates, followed by implication introduction and modus ponens, yields \(O\).  The actual codes disappear.

\begin{theorem}[Stable-observer erasure]
\label{thm:stable-erasure}
Let \(R_i(y_i)\) be finite producer relations and suppose object-iPRA proves:
\begin{align*}
 \Gamma&\vdash\exists y_i<B_i\ R_i(y_i)\qquad(1\le i\le k),\\
 R_1(y_1),\ldots,R_k(y_k),\Gamma&\vdash O
\end{align*}
with the conclusion independent of the bound variables \(y_i\).  Then object-iPRA proves \(\Gamma\vdash O\).
\end{theorem}

\begin{proof}
Apply existential elimination successively.  Freshness is satisfied because \(O\) and \(\Gamma\) do not contain the witness variables.  Equivalently, if exact certificate terms are already available, use the repository's \texttt{observableFromLocalStability} combinator, which implements the same proof by implication introduction and elimination \cite{repo-observable}.
\end{proof}

This theorem should be preferred to least selection whenever the consumer is genuinely witness-independent.  It produces smaller proof terms and avoids artificial canonicity claims.

\subsection{Metric-stable observers}

A terminal observable may not be exactly equal for all witnesses but may vary by a controlled amount.  Let \(O(y)\) be a rational observation and suppose all viable witnesses at precision \(p\) satisfy
\[
 d(O(a),O(b))\le\eta_p.
\]
Then any least-witness output is canonical up to \(\eta_p\).  If \(\eta_p\to0\) with an explicit modulus, the outputs define the same represented limit.  This is the approximate analogue of stable erasure.

In the object theory, replace equality by a scaled rational inequality.  The proof is again finite: produce any two valid certificates and apply the local stability estimate.  Extensional equality of represented objects is then defined observationally rather than as equality of codes.

\subsection{Root collisions}

Consider a monic polynomial with approximate coefficients.  At precision \(p\), a certified root producer may output a multiset of rational disks whose union contains all roots with multiplicity.  Individual disk labels can permute or merge as coefficients vary.  The correct stable observables include:

\begin{itemize}
\item the multiset of disks modulo permutation;
\item symmetric polynomials of root approximants;
\item clusters separated from the remaining spectrum;
\item root counts in regions whose boundary has a certified gap.
\end{itemize}

An ordered root name is legitimate only under a separation certificate showing that distinct clusters remain at least \(\gamma>0\) apart and the approximation error is less than, say, \(\gamma/3\).  Then nearest-cluster matching is unique and primitive recursive.

\subsection{Spectral data}

The same pattern applies to matrices.  Ordered eigenvalues of a self-adjoint matrix are stable as a multiset and can be named using variational or characteristic-polynomial data.  Eigenvectors are not stable at repeated eigenvalues.  Stable outputs are:

\begin{itemize}
\item the characteristic polynomial;
\item the ordered eigenvalue list with multiplicity;
\item spectral clusters;
\item a projector onto a cluster separated by a positive gap;
\item polynomial functions of the matrix.
\end{itemize}

A projector requires a gap guard.  An individual eigenvector additionally requires a one-dimensionality certificate and a sign convention separated from zero in a chosen coordinate.  The framework thereby turns an apparent logical obstruction into a precise interface condition.

\subsection{Metastability versus a Cauchy modulus}

A metastability statement has the form
\[
 \forall p\,\forall g\,\exists N\le\Phi(p,g)\,
 \forall i,j\in[N,N+g(N)]\ d(x_i,x_j)\le2^{-p}.
\]
It gives arbitrarily long stable finite windows but not, in general, a computable tail modulus.  PFIGs can represent and prove the metastable theorem as a finite search contract: for fixed \(p,g\), search \(N\le\Phi(p,g)\) and check the bounded window.  They must not then advertise \((x_n)\) as a regular Cauchy name unless an additional tail modulus is supplied.

Shrinking viability is stronger because its diameter and motion hypotheses directly give a summable tail.  The framework records the two outputs as different interface types:

\begin{itemize}
\item \texttt{MetastableSequence}: a finite-window search contract;
\item \texttt{RegularName}: a primitive-recursive approximant with a tail modulus.
\end{itemize}

\subsection{Formalization shape}

The Lean development should prove three layers.

\begin{enumerate}
\item A semantic theorem on finite lists or bounded predicates: viability, predecessor extraction, and the adjacent coherence inequality.
\item An object-iPRA theorem reflecting the viability table and least-pair selection.
\item A rational metric library proving finite triangle sums and converting the tail certificate into the regular-name formula.
\end{enumerate}

The first layer can be executable and tested independently.  The second and third are the proof-producing bridge.  The theorem ledger must not mark the regular-name result complete until the object reflection, not just the semantic inequality, has been constructed.

\section{Guarded operations, extensionality, and discontinuous choices}
\label{sec:guarded}

An effective presentation is useful only if its operations respect the represented object rather than accidental code choices.  At the same time, many familiar operations are partial or discontinuous.  This section gives a uniform treatment: totalize the numerical program, state the mathematical law under an explicit guard, and formulate equality observationally.

\subsection{Total programs with guarded laws}

A primitive-recursive function is total on natural-number codes.  Division, normalization, chart inversion, pivoting, logarithms, and branch selection are therefore implemented with a default output on invalid inputs.  Correctness is stated conditionally.

\begin{definition}[Guarded operation]
A guarded operation from presentation \(X\) to presentation \(Y\) consists of:

\begin{enumerate}
\item a total primitive-recursive code function \(F\);
\item a quantifier-free or bounded-decidable guard \(G(g,x)\), where \(g\) is quantitative margin data;
\item a demand function \(d(p,g,x)\);
\item an object proof that, under \(G(g,x)\), the output observations of \(F(x)\) satisfy the intended law at precision \(p\).
\end{enumerate}
\end{definition}

The default branch is never used in the guarded theorem.  This keeps the program total without pretending that the mathematical operation is total.

\subsection{Apartness is data}

For a represented real \(x\), the proposition \(x\neq0\) is negative and may not reveal how accurately one must approximate \(x\) before its sign or reciprocal is determined.  A usable inverse interface consumes an integer \(g\) such that
\[
 |x|\ge 2^{-g}.
\]
If \(r_q\) approximates \(x\) within \(2^{-q}\), choosing \(q\ge g+2\) guarantees \(|r_q|\ge 2^{-g-1}\).  Rational inversion is then safe, and a direct estimate gives a demand for the desired output precision.

The same principle applies to:

\begin{itemize}
\item a pivot: supply a lower bound on the pivot magnitude;
\item a logarithm: supply a positive lower bound on the argument;
\item a normalized vector: supply a lower bound on its norm;
\item an inverse chart: supply a distance from the chart boundary;
\item a spectral projector: supply a gap separating the cluster;
\item a simple root: supply a derivative or separation bound.
\end{itemize}

A theorem that assumes only a nonzero proposition should either derive such a bound from stronger exact input data or remain a conditional relation rather than a computable operation.

\subsection{Branch stabilization under a gap}

Let alternatives be labeled by a finite set \(L\).  Suppose an exact score \(s_\ell(x)\) is represented by rational approximations \(\widehat s_\ell(q,x)\) with error at most \(e_q\).  Assume one label \(\ell_0\) satisfies
\[
 s_{\ell_0}(x)+\gamma\le s_\ell(x)
 \qquad(\ell\ne\ell_0)
\]
for a supplied rational gap \(\gamma>0\).

\begin{theorem}[Guarded branch stabilization]
\label{thm:branch-stability}
If \(2e_q<\gamma\), then the least approximate minimizer of \(\widehat s_\ell(q,x)\) is \(\ell_0\).  Hence the label is primitive-recursively selectable from sufficiently precise observations and the gap certificate.
\end{theorem}

\begin{proof}
For \(\ell\ne\ell_0\),
\[
 \widehat s_{\ell_0}
 \le s_{\ell_0}+e_q
 \le s_\ell-\gamma+e_q
 < s_\ell-e_q
 \le \widehat s_\ell.
\]
Thus \(\ell_0\) is the unique approximate minimizer.  Every comparison is rational and bounded.  The object proof is a chain of coded integer inequalities.
\end{proof}

This theorem is the generic resolution for exact labels under separation.  Without \(\gamma\), the appropriate output is a finite set of near-minimizers or a stable aggregate.

\subsection{Observational setoids}

Codes are rarely canonical.  Define equality at precision \(p\) by a bounded relation
\[
 x\approx_p y.
\]
For exact algebraic codes this may be literal equality or an ideal-membership certificate.  For rational names it is metric closeness.  For finite covers it is local equivalence on every chart.  For root clouds it is a finite matching relation or equality of symmetric observations.

A presentation should prove:

\begin{enumerate}
\item reflexivity at every precision;
\item symmetry;
\item a quantitative transitivity law, often
\(x\approx_{p+c}y\wedge y\approx_{p+c}z\to x\approx_p z\);
\item restriction monotonicity: higher-precision equivalence implies lower-precision equivalence;
\item extensionality of every operation with an explicit demand map.
\end{enumerate}

These are finite contracts and can be compiled by PFIGs.  No quotient type is needed inside the object theory.  The quotient is a semantic interpretation of the observational setoid.

\subsection{Relational extensionality}

For a finite relation \(R(x,a,p)\), code equality of answer clouds is too strong.  The right extensionality notion is finite back-and-forth:
\begin{align*}
 x\approx_{d(p)}x'\ \wedge\ R(x,a,d(p))
 &\to \exists a'<B'\bigl(R(x',a',p)\wedge a\sim_p a'\bigr),\\
 x\approx_{d(p)}x'\ \wedge\ R(x',a',d(p))
 &\to \exists a<B\bigl(R(x,a,p)\wedge a\sim_p a'\bigr).
\end{align*}
All quantifiers are bounded.  This says that equivalent inputs produce observationally equivalent answer clouds, even if their codes, cardinalities, or orderings differ.

When a stable observer is invariant under \(\sim_p\), this back-and-forth is enough to prove extensionality of the observed output.

\subsection{Permutation and symmetry}

Many finite outputs are unordered collections.  Two practical representations are available.

\paragraph{Canonical sorting.}
If codes have a decidable total order that respects the desired identity, sort them and compare sorted lists.  This is suitable for exact rational factors or intervals with separated centers.

\paragraph{Matching certificates.}
Represent equivalence by a bounded permutation or bipartite matching certificate.  This is suitable near collisions, where canonical order can be unstable.  Hall-type conditions are finite and Booleanizable; in small fixed dimensions a direct permutation certificate is simplest.

The framework does not force one representation.  It requires the operation's extensionality proof to match the chosen representation.

\subsection{Projective and quotient-like outputs}

A direction, eigenvector, homogeneous coordinate, or rotation axis is naturally defined only up to scaling or sign.  Instead of imposing a discontinuous normalization, represent it projectively:
\[
 v\sim w\quad\Longleftrightarrow\quad
 \exists \lambda\ne0\; w=\lambda v,
\]
or, in finite exact arithmetic, by vanishing of all \(2\times2\) minors \(v_iw_j-v_jw_i\).  The latter is quantifier-free.  Approximate projective equivalence is expressed by small minors and positive norm guards.

A chart may be selected only when a coordinate has a positive lower bound.  Otherwise, retain a finite relation of admissible charts.  Stable projective observables, such as the rank-one projector \(vv^*/\|v\|^2\), can be produced when the norm and spectral gap are guarded.

\subsection{Rotation axes}

At the identity rotation, every unit axis is valid.  A total exact axis selector would have to make a discontinuous choice.  The constructive output should instead be one of:

\begin{itemize}
\item a unit quaternion modulo sign;
\item an antisymmetric infinitesimal or projective axis relation;
\item a finite cloud of approximate axes;
\item an exact axis under a positive distance-from-identity guard.
\end{itemize}

Rodrigues' formula itself is a polynomial/rational identity under a suitable parameterization and can be checked finitely.  The obstruction concerns the selector, not the identity.

\subsection{Piecewise definitions}

A piecewise operation with decidable exact branches is compiled by Boolean case analysis.  A piecewise operation whose branch depends on the sign of a represented real must use one of three interfaces:

\begin{enumerate}
\item \textbf{guarded branch}: a sign-separation bound is an input;
\item \textbf{relational branch}: both branches are returned when zero cannot be excluded;
\item \textbf{stable merge}: prove that both branches give the same observable near the boundary.
\end{enumerate}

This trichotomy should be recorded explicitly in theorem statements.  It prevents hidden uses of limited omniscience principles.

\subsection{Extensional composition}

\begin{theorem}[Composition of precision contracts]
\label{thm:precision-composition}
Suppose \(F:X\to Y\) is extensional with demand \(d_F\), and \(G:Y\to Z\) is extensional with demand \(d_G\).  Then \(G\circ F\) is extensional with demand
\[
 d_{G\circ F}(p,x)=d_F(d_G(p,Fx),x),
\]
or any primitive-recursive upper bound on this expression.
\end{theorem}

\begin{proof}
If inputs agree at the composite demand, extensionality of \(F\) gives agreement of outputs at the precision requested by \(G\).  Extensionality of \(G\) then gives agreement at terminal precision \(p\).  In a code-dependent demand, the expression may refer to an approximate output of \(F\); the graph's backward pass computes a safe majorant using the guard and finite observations.  The object proof is substitution followed by cut.
\end{proof}

For relational functions, replace direct agreement by finite back-and-forth and compose the bounded witnesses.  For metric relations, add the error budgets and use a stronger internal precision.

\subsection{The role of defaults}

Defaults should be semantically inert.  Every theorem about a totalized operation has the form
\[
 G(g,x)\to \Post(F(g,x),x).
\]
No theorem should expose the default as a meaningful output when \(G\) fails.  If downstream nodes need to continue under guard failure, the output type should include an explicit status bit and the terminal statement should branch on that bit.  This makes failure observable and prevents accidental use of undefined data.

\subsection{Resolution achieved}

The guarded-extensional layer resolves several recurrent problems:

\begin{itemize}
\item code equality is replaced by observational equality;
\item quotients are represented by certified setoids rather than object-theoretic quotient types;
\item discontinuous labels become relations or guarded exact outputs;
\item partial operations become total programs with conditional laws;
\item precision bookkeeping composes through explicit demand maps.
\end{itemize}

The remaining obligation is mathematical: produce the relevant gap, apartness, coverage, or continuity modulus.  The framework makes this obligation impossible to hide.

\section{Finite presentations, covers, and local-to-global gluing}
\label{sec:gluing}

Many advanced mathematical concepts appear to require sets, quotient objects, sheaves, manifolds, or categories.  In the intended applications, however, a proof uses only finitely many charts, generators, relations, overlaps, and observations at a fixed request.  PFIGs capture this finite-use content directly.

\subsection{Certified finite presentations}

A certified finite presentation consists of:

\begin{enumerate}
\item a natural-number code for finite generators or local pieces;
\item primitive-recursive decoding functions and validity predicates;
\item bounded relations expressing equations, incidences, or overlaps;
\item finite certificates whose checkers reflect the desired mathematical laws;
\item observation and restriction functions at requested precision.
\end{enumerate}

The presentation need not be canonical.  Its equivalence relation is part of the interface.  A proof may use a tuple of representatives throughout and invoke extensionality only when comparing outputs.

\subsection{Finite cover data}

Let a finite cover have indices \(i<C\).  A local object is coded by \(s_i\), and an overlap between \(i\) and \(j\) has a bounded compatibility predicate
\[
  \mathsf{Compat}(i,j,s_i,s_j,w_{ij}),
\]
where \(w_{ij}\) is an optional certificate.  The global code is the tuple
\[
  S=\langle s_0,\ldots,s_{C-1},w_{00},w_{01},\ldots\rangle.
\]
Validity is a bounded conjunction over all indices and overlaps, hence Booleanizable.

A local operation \(F_i\) acts on each \(s_i\).  To define a global operation, prove an overlap theorem:
\[
 \mathsf{Compat}(s_i,s_j)	o
 \mathsf{Compat}(F_i(s_i),F_j(s_j)).
\]
A bounded loop applies \(F_i\) to all pieces and checks all new overlaps.

\subsection{Finite gluing theorem}

\begin{theorem}[Certified finite gluing]
\label{thm:finite-gluing}
Suppose a fixed finite cover has:

\begin{enumerate}
\item primitive-recursive local code bounds and validity predicates;
\item primitive-recursive overlap restrictions;
\item a bounded compatibility checker with an object reflection theorem;
\item a local construction on every chart with an object correctness proof;
\item a proof that local outputs agree on every overlap;
\item a finite uniqueness or observational extensionality theorem.
\end{enumerate}

Then PFIG compilation produces a global compatible tuple and an object-iPRA derivation of its global specification.  If a finite bounded set of global representatives is provided, least search can additionally produce a canonical representative; otherwise the compatible tuple itself is the global code.
\end{theorem}

\begin{proof}
Run the local construction over \(i<C\) by bounded recursion, storing each output.  The loop invariant says that every completed local output is valid and satisfies its local specification.  Next run a bounded double loop over pairs \((i,j)\), checking overlap compatibility and reflecting each true Boolean to the overlap formula.  The supplied local-to-overlap theorem proves that every check succeeds.

The resulting tuple is a global code by definition.  A global observer is evaluated componentwise and shown independent of equivalent tuples using the finite uniqueness/extensionality theorem.  If a separately bounded list of global representatives exists, search it for the least representative equivalent to the tuple.  No quotient construction or unbounded choice is needed.
\end{proof}

\subsection{Why tuples are often better than quotients}

Inside weak arithmetic, quotient types create unnecessary proof obligations.  A compatible tuple already supports all finite observations required by downstream nodes.  Equality of global objects is defined by componentwise observational equivalence plus overlap compatibility.  Operations are proved extensional with respect to this relation.

A quotient may still be used in the external semantics to explain the mathematics.  The object theory manipulates representatives and certificates.  This is analogous to representing rational numbers by integer pairs and proving cross-multiplication congruence instead of forming a quotient inside arithmetic.

\subsection{Algebraic presentations}

A finitely presented algebraic object can be encoded by generators, relations, and bounded expression trees.  Equality may be certified by a relation witness:

\begin{itemize}
\item an ideal-combination certificate for polynomial quotient rings;
\item a word-reduction trace for finitely presented groups in decidable fragments;
\item a row-operation certificate for modules;
\item a basis-change matrix and inverse certificate for finite-dimensional spaces.
\end{itemize}

A checker verifies the certificate by finite arithmetic.  Composition of maps is substitution of expression trees followed by bounded normalization or direct certificate checking.  A PFIG node need not decide all equalities in the presented structure; it only verifies the finite equalities used by the proof.

\subsection{Localization and denominators}

Localization requires careful handling of zero divisors.  Rather than identify fractions solely by cross multiplication, use the appropriate finite certificate, for example a killing exponent \(n\) with
\[
 s^n(t a-s b)=0
\]
for denominators generated by \(s,t\).  The exponent and algebraic identity are finite data.  Operations on fraction codes are total, while correctness is conditioned on denominator validity and the equality certificate.

When inversion of a represented analytic quantity is needed, use the guarded apartness interface of \cref{sec:guarded}; algebraic localization certificates and analytic lower bounds are different kinds of evidence and should not be conflated.

\subsection{Finite categories and diagrams}

A finite category is represented by:

\begin{itemize}
\item a bound on objects and morphisms;
\item source, target, identity, and composition tables;
\item Boolean checks of the category axioms over bounded indices.
\end{itemize}

A finite diagram is a bounded table of object and morphism codes satisfying functoriality checks.  Finite limits, colimits, natural transformations, and adjunction data can be represented by candidate cones/cocones and universal-property certificates over explicitly bounded hom-sets.

For an arbitrary fixed program family, categorical composition is a PFIG composition theorem.  What is not available is an internal category of all primitive-recursive programs with a total universal evaluator.  The fixed finite presentation boundary remains essential.

\subsection{Sheaf-like gluing}

For a finite cover, a presheaf is represented by finite section codes and restriction maps.  The sheaf condition used by a concrete proof is:

\begin{enumerate}
\item a bounded compatibility check for a finite family of local sections;
\item a primitive-recursive gluing constructor or a bounded relation of possible gluings;
\item a reflection proof that the output restricts to each local section;
\item a uniqueness theorem up to observational equality.
\end{enumerate}

This is exactly \cref{thm:finite-gluing}.  Infinite topological quantification is replaced by the finite cover actually used at the requested stage.  For compact metric presentations, a finite net or finite subcover certificate is supplied by a producer node.

\subsection{Simplicial and cubical observation towers}

Topology is represented through finite complexes at increasing resolution.  At precision \(p\), a space has a finite simplicial or cubical complex \(K_p\).  A map is represented by a finite carrier relation: each vertex maps to a nonempty finite cloud of target vertices lying in a certified carrier simplex.

Any two selections from the same carrier are contiguous; an explicit prism operator gives a chain homotopy.  Thus the induced homology map is independent of tie-breaking.  This is a paradigmatic stable observer:

\begin{itemize}
\item internal vertex choices are relational;
\item chain maps may differ at code level;
\item the induced homology observable is stable under contiguity.
\end{itemize}

Finite boundary matrices, cycles, boundaries, and homology certificates are exact bounded algebra.  A theorem such as a Sperner parity argument becomes a Boolean certificate over a finite grid.

\subsection{Local solutions and ODE patches}

A local analytic construction may produce rational polynomial or interval enclosures on short time steps.  A global finite-time solution is a compatible list of patches.  Each patch has:

\begin{itemize}
\item an initial box;
\item a polynomial or rational approximation;
\item a residual bound;
\item a Lipschitz or contraction certificate;
\item an endpoint enclosure compatible with the next patch.
\end{itemize}

A bounded loop glues the patches.  The terminal result is valid on a supplied finite time interval.  No unbounded existence theorem is hidden: the step size, number of patches, and error budget are explicit.

\subsection{Local-to-global descent}

A finite descent argument has states \(s\), a rank \(r(s)\), local reduction steps, and a terminal normal-form relation.  If every nonterminal step decreases a rank bounded by a primitive-recursive function, totalize the update and run it for that bound.  The invariant records preservation of the mathematical class and rank decrease.  Terminal extraction yields the global conclusion.

This pattern covers Euclidean algorithms, reduction of quadratic forms, finite descent on explicit covering data, and normalization of finite presentations.  It does not supply missing rational points on covering varieties; those remain producer obligations.

\subsection{Complete finite lists from bounds}

The most important local-to-global instance for Diophantine geometry is bounded enumeration.

\begin{theorem}[Height bound implies complete list]
\label{thm:height-to-list}
Let \(P(x)\) be a decidable predicate on coded objects and let \(h(x)\) be a primitive-recursive height.  Suppose there is a primitive-recursive enumeration of all valid codes of height at most \(B\).  Then from a bound
\[
 P(x)\to h(x)\le B
\]
one can primitive-recursively produce a finite list containing exactly all \(P\)-objects up to code equivalence, together with an object-iPRA proof of completeness.
\end{theorem}

\begin{proof}
Enumerate every code in the finite height box, test validity and \(P\), and retain the least representative of each decidable equivalence class.  Soundness follows from the checker.  For completeness, any \(P(x)\) has height at most \(B\), hence appears in the finite enumeration; duplicate elimination preserves an equivalent representative.  All loops and quantifiers are bounded.
\end{proof}

\begin{theorem}[Complete list implies height bound]
\label{thm:list-to-height}
If a supplied finite list is sound and complete for \(P\), then the maximum height of its members is a bound for every \(P\)-object.
\end{theorem}

\begin{proof}
Compute the finite maximum.  Completeness gives, for every \(P(x)\), an equivalent list member.  Height invariance or a bounded comparison under equivalence yields the claim.
\end{proof}

Together, the two theorems show that, for decidable coded objects with finite height boxes, an explicit height bound and a complete finite list are computationally equivalent data.  Mere finiteness or a cardinality bound supplies neither.

\subsection{Application to difficult Diophantine theorems}

For Siegel- or Faltings-type statements, the PFIG framework completely resolves the verifier side:

\begin{itemize}
\item check a proposed height bound;
\item enumerate the bounded height box;
\item verify that each listed point lies on the variety;
\item prove completeness from the bound;
\item compute any finite observable of the point set.
\end{itemize}

The unresolved mathematical producer is an explicit primitive-recursive bound uniform in the intended input family.  A proof of abstract finiteness does not fill this slot.  Special families with an effective descent or Runge bound fit the framework immediately.

\subsection{Summary}

Finite presentations make sophisticated structures compatible with iPRA because every actual proof request uses only finitely many codes and checks.  The gluing theorem, observational setoids, bounded search, and trace reflection supply a general local-to-global mechanism.  The framework is not a claim that all global existence theorems are effective; it is a precise account of the finite data that would make them effective.

\section{Classical-to-intuitionistic transfer for the retained fragment}
\label{sec:negative-transfer}

Some repository proofs are easiest to discover classically, even though their final arithmetic statement is stable and constructive.  A disciplined negative translation can recover such proofs in iPRA without adding excluded middle to the target theory.  The existing \texttt{NegativeSyntax.lean} file defines the translation, proves its substitution laws, translates every arithmetic axiom and excluded-middle instance, and records the remaining proof-translation and stability statements as goals.  This section supplies the full mathematical argument and an exact Lean completion path.

\subsection{The translation}

Write \(\neg A\) for \(A\to\bot\).  Define an atom-preserving translation \(A\mapsto A^N\) by
\begin{align*}
 (s=t)^N&=(s=t),& \bot^N&=\bot,\\
 (A\wedge B)^N&=A^N\wedge B^N,&
 (A\to B)^N&=A^N\to B^N,\\
 (A\vee B)^N&=\neg(\neg A^N\wedge\neg B^N),&
 (\forall x A)^N&=\forall x A^N,\\
 (\exists x A)^N&=\neg\forall x\neg A^N.
\end{align*}
This is precisely the translation encoded in the repository file.  It commutes with term substitution and preserves quantifier-freeness; those syntax lemmas are already proved there.

Let \(\mathsf{CPRAAx}\) be the source theory consisting of the ordinary iPRA arithmetic axioms plus every excluded-middle instance \(A\vee\neg A\).  The target remains \(\PRAAx\), with intuitionistic logic only.

\subsection{Decidability and stability of numerical equality}

\begin{lemma}[Equality decidability]
\label{lem:eq-decidable}
For all primitive-recursive terms \(s,t\), iPRA proves
\[
  s=t\ \vee\ \neg(s=t).
\]
\end{lemma}

\begin{proof}
Use the primitive-recursive Boolean equality function \(\mathsf{eq}(s,t)\).  Bounded Booleanization gives
\(
 [\mathsf{eq}(s,t)]\leftrightarrow s=t
\)
and proves that the Boolean is either \(0\) or \(1\).  A finite Boolean case split yields the disjunction.
\end{proof}

\begin{corollary}[Equality stability]
\label{cor:eq-stability}
For all terms \(s,t\), iPRA proves
\[
  \neg\neg(s=t)\to s=t.
\]
\end{corollary}

\begin{proof}
Case-split using \cref{lem:eq-decidable}.  Equality is immediate in the positive case; in the negative case, the double negation produces a contradiction.
\end{proof}

This is the missing arithmetic fact that permits an atom-preserving negative translation.  No global excluded middle is used; only the decidability of natural-number equality, proved by a fixed Boolean computation.

\subsection{Stability of every translated formula}

\begin{theorem}[Negative stability]
\label{thm:negative-stability}
For every formula \(A\), iPRA proves
\[
  \neg\neg A^N\to A^N.
\]
\end{theorem}

\begin{proof}
Proceed by structural induction on \(A\).

\paragraph{Atomic equality.}
Use \cref{cor:eq-stability}.  For \(\bot\), note that \(\bot\to\bot\) is provable; applying \(\neg\neg\bot\) to it yields \(\bot\).

\paragraph{Conjunction.}
Assume \(\neg\neg(A^N\wedge B^N)\).  To obtain \(A^N\), by the induction hypothesis it suffices to prove \(\neg\neg A^N\).  If \(\neg A^N\), then \(\neg(A^N\wedge B^N)\), contradicting the assumption.  Similarly obtain \(B^N\), then introduce the conjunction.

\paragraph{Disjunction.}
The translation \((A\vee B)^N\) is itself a negation.  Every negation is intuitionistically stable: from \(\neg\neg\neg C\) and \(C\), derive \(\neg\neg C\), then contradiction; hence \(\neg C\).

\paragraph{Implication.}
Assume \(\neg\neg(A^N\to B^N)\) and suppose \(A^N\).  By stability of \(B^N\), it suffices to prove \(\neg\neg B^N\).  If \(\neg B^N\), then under the fixed assumption \(A^N\), every function \(A^N\to B^N\) yields a contradiction.  Thus \(\neg(A^N\to B^N)\), contradicting the double negation.  Hence \(B^N\).

\paragraph{Universal quantification.}
Assume \(\neg\neg\forall x A^N(x)\).  Fix arbitrary \(x\).  By the induction hypothesis, it suffices to prove \(\neg\neg A^N(x)\).  If \(\neg A^N(x)\), then any proof of \(\forall x A^N(x)\) yields a contradiction by specialization, contrary to the double negation.  Generalize over \(x\).

\paragraph{Existential quantification.}
The translation \((\exists x A)^N=\neg\forall x\neg A^N\) is a negation, hence stable as in the disjunction case.
\end{proof}

The proof uses only intuitionistic logic plus equality stability at atomic leaves.

\subsection{Translation of derivations}

\begin{theorem}[Proof translation]
\label{thm:proof-translation}
If
\[
  \Gamma\vdash_{\mathsf{CPRAAx}} A,
\]
then
\[
  \Gamma^N\vdash_{\PRAAx} A^N,
\]
where \(\Gamma^N\) is the pointwise translation of the context.
\end{theorem}

\begin{proof}
Induct on the source derivation.  The hypothesis, conjunction, implication, universal, equality-reflexivity, and equality-substitution cases follow directly from the corresponding target rule, the induction hypotheses, and commutation of translation with substitution.  Explosion is immediate because \(\bot^N=\bot\).  Arithmetic axioms and excluded middle have already translated target derivations in the repository file: arithmetic equations are unchanged or translated structurally, and
\[
 (A\vee\neg A)^N
 =\neg(\neg A^N\wedge\neg(A^N\to\bot))
\]
after simplification is intuitionistically derivable.

It remains to spell out the nontrivial disjunction and existential rules.

\paragraph{Disjunction introduction.}
Suppose the induction hypothesis gives \(A^N\).  To prove
\(\neg(\neg A^N\wedge\neg B^N)\), assume the conjunction and apply its first component to \(A^N\).  The right introduction case is symmetric.

\paragraph{Disjunction elimination.}
Suppose we have
\begin{align*}
 h&:\neg(\neg A^N\wedge\neg B^N),\\
 f&:A^N\to C^N,\\
 g&:B^N\to C^N.
\end{align*}
By \cref{thm:negative-stability}, it suffices to prove \(\neg\neg C^N\).  Assume \(k:\neg C^N\).  Then
\[
  \lambda a.\,k(f(a)):\neg A^N,
  \qquad
  \lambda b.\,k(g(b)):\neg B^N.
\]
Their conjunction contradicts \(h\).  Therefore \(\neg\neg C^N\), and stability yields \(C^N\).  In the natural-deduction formalization, \(f\) and \(g\) are the translated branch derivations under extended contexts; implication introduction internalizes them before the argument above.

\paragraph{Existential introduction.}
From \(A^N(t)\), prove \(\neg\forall x\neg A^N(x)\): assume the universal negation, specialize at \(t\), and obtain contradiction.

\paragraph{Existential elimination.}
Suppose
\[
 h:\neg\forall x\neg A^N(x)
\]
after translating the major premise, and suppose the translated branch proves \(C^N\) from a fresh assumption \(A^N(x)\).  By stability it suffices to prove \(\neg\neg C^N\).  Assume \(k:\neg C^N\).  For arbitrary \(x\), the branch derivation and \(k\) show \(\neg A^N(x)\).  Universal introduction gives \(\forall x\neg A^N(x)\), contradicting \(h\).  Hence \(\neg\neg C^N\), and stability gives \(C^N\).  The freshness condition is exactly the one already enforced by the source and target existential-elimination constructors.

\paragraph{Equality substitution.}
The translated equality premises are the same equalities.  Apply target equality substitution to the formula \(A^N\).  The syntax lemma
\((A[\sigma])^N=A^N[\sigma]\)
identifies the result with the induction-hypothesis conclusion.

Thus every source constructor has a target derivation, completing the induction.
\end{proof}

\subsection{Quantifier-free equivalence}

The translation is not generally equivalent to the original formula intuitionistically.  It is equivalent on decidable formulas.

\begin{theorem}[Quantifier-free recovery]
\label{thm:qf-recovery}
For every quantifier-free formula \(A\), iPRA proves
\[
  A^N\leftrightarrow A.
\]
\end{theorem}

\begin{proof}
By \cref{thm:booleanization}, every fixed quantifier-free formula is decidable in iPRA.  One may prove the result structurally.

Atomic equalities, \(\bot\), conjunctions, and implications follow directly from the induction hypotheses.  For disjunction, the forward direction is the only interesting one.  Suppose
\(\neg(\neg A^N\wedge\neg B^N)\).  Decide \(A\).  If \(A\), take the left disjunct.  Otherwise decide \(B\).  If \(B\), take the right disjunct.  If neither holds, the induction equivalences give \(\neg A^N\) and \(\neg B^N\), contradicting the translated disjunction.  The reverse direction is disjunction introduction as in the proof-translation theorem.
\end{proof}

The use of formula-specific Booleanization is legitimate because the formal theorem is proved by recursion on formula syntax, generating the corresponding Boolean term and reflection proof.

\subsection{The retained class}

Define \(\mathsf{Retained}\) to be the least class containing all quantifier-free formulas and closed under \(\wedge\), \(\to\), and \(\forall\).  This is the class already specified in the repository blueprint.

\begin{theorem}[Retained equivalence]
\label{thm:retained-equivalence}
If \(A\) is retained, then iPRA proves
\[
  A^N\leftrightarrow A.
\]
\end{theorem}

\begin{proof}
Induct on the retained derivation.  The quantifier-free case is \cref{thm:qf-recovery}.  Conjunction and implication use the induction equivalences and ordinary intuitionistic reasoning.  For universal quantification, use the pointwise equivalence under an arbitrary bound variable and introduce the universal in each direction.
\end{proof}

The class may be enlarged whenever another formula family is proved stable and equivalent to its translation.  For example, bounded formulas are retained after Booleanization, and certain existential statements with primitive-recursive unique witnesses can be rewritten as quantifier-free graph assertions.

\subsection{Retained transfer theorem}

\begin{theorem}[Classical-to-intuitionistic retained transfer]
\label{thm:retained-transfer}
Assume every formula in \(\Gamma\) and the conclusion \(A\) is retained.  Then
\[
 \Gamma\vdash_{\mathsf{CPRAAx}}A
 \quad\Longrightarrow\quad
 \Gamma\vdash_{\PRAAx}A.
\]
Conversely, every iPRA derivation lifts trivially to the classical source theory, so derivability agrees on retained sequents.
\end{theorem}

\begin{proof}
Translate the classical derivation using \cref{thm:proof-translation}:
\[
 \Gamma^N\vdash_{\PRAAx}A^N.
\]
For every assumption \(B\in\Gamma\), retained equivalence gives a derivation \(B\vdash B^N\).  Apply simultaneous cut to replace the translated context by \(\Gamma\), obtaining \(\Gamma\vdash A^N\).  Retained equivalence for the conclusion gives \(A^N\to A\); modus ponens yields \(A\).

For the converse, map each \(\PRAAx\) axiom constructor to the arithmetic constructor of \(\mathsf{CPRAAx}\) and recurse over the unchanged intuitionistic logical derivation.
\end{proof}

\subsection{Interaction with PFIGs}

The transfer theorem is useful in two ways.

\begin{enumerate}
\item A finite arithmetic lemma may be proved classically in a source development and transferred automatically when its assumptions and conclusion are retained.
\item A PFIG terminal theorem often has a retained shape: universal closure of implications between bounded or quantifier-free certificate assertions.  Classical finite combinatorics can therefore be used as a proof-discovery layer, while the final result is recovered in iPRA.
\end{enumerate}

This does not make arbitrary classical mathematics available in iPRA.  Conclusions with essential positive disjunctions or existentials may not be retained.  For such theorems, the PFIG must expose explicit witnesses or finite relations.

\subsection{Lean completion plan for the existing goals}

The repository's \texttt{NegativeSyntax.lean} already contains the syntax translation, substitution lemmas, translated induction axioms, translated excluded middle, and translated arithmetic axioms.  The remaining implementation can be organized as follows.

\begin{enumerate}
\item \texttt{QFDecidable.lean}: construct a PR Boolean for every QF formula and prove its object reflection.  Derive \texttt{equalityStability}.
\item \texttt{NegativeStability.lean}: recurse on \(\Fm\) and implement \cref{thm:negative-stability}.  Package the standard triple-negation lemma and stable implication/universal combinators.
\item \path{NegativeProof.lean}: recurse on \path{Prf ClassicalPRAAx}.  The \texttt{orE} and \texttt{exE} branches use \path{negativeStability}; all substitution equalities are discharged by the existing \texttt{negative\_subst} lemma.
\item \texttt{Retained.lean}: prove QF recovery by structural recursion using decidability, then retained equivalence.
\item \texttt{RetainedTransfer.lean}: use \texttt{Prf.trans} to replace translated assumptions and prove the source-theory lifting direction.
\end{enumerate}

The difficult part is de Bruijn bookkeeping in \texttt{exE} and \texttt{allI}, not a missing mathematical principle.  The theorem statements in the accompanying Lean blueprint match these milestones.

\subsection{Status of the difficulty}

At the human mathematical level, the five open goals named by the repository file are resolved by the arguments above.  At the formal level, they remain to be entered and kernel-checked.  The important point is that no additional axiom, Markov principle, or excluded-middle rule is required: numerical equality decidability and translated-formula stability suffice.

\section{Reusable embedding patterns across mathematical domains}
\label{sec:domains}

The purpose of a general framework is not to replace theorem-specific mathematics but to make its reusable computational and logical skeleton explicit.  This section records standard PFIG patterns for the main domains represented in the repository.

\subsection{Finite algebra and polynomial identities}

\paragraph{Data.}
Integers are signed pairs or a canonical sign-and-magnitude code.  Rationals are integer pairs with positive denominators; reduction to lowest terms is optional because equality and order are checked by cross multiplication.  Vectors, matrices, and polynomials are bounded tables.

\paragraph{Operations.}
Addition, multiplication, evaluation, differentiation, coefficient extraction, finite sums, determinants, resultants, and substitutions are primitive recursive.  Their algebraic identities are proved by equality induction or reflected certificates.

\paragraph{PFIG pattern.}
An exact transducer builds the expression, a bounded loop evaluates table entries, and a certificate checker verifies the final identity.  If a symbolic computation is imported, provide an expansion or row-operation trace.  The terminal theorem is stable under harmless representation choices such as unreduced fractions.

\paragraph{Difficult cases.}
Factorization and root extraction are relational.  Return a finite factor list modulo permutation, or root clusters with multiplicity certificates.  A canonical ordered list requires a decidable stable ordering or a separation bound.

\subsection{Linear algebra and spectral theory}

Matrices of variable finite dimension are coded by a dimension and a bounded entry table.  Matrix multiplication and characteristic polynomials are finite folds.  Gaussian elimination is totalized: pivot search is bounded, and a singularity/status bit is carried through the state.

A proof of rank, determinant, inverse, or solution existence should use a certificate:

\begin{itemize}
\item a nonzero minor for a rank lower bound;
\item a row-reduction trace for rank upper bounds and consistency;
\item an adjugate identity for inversion under nonzero determinant;
\item a kernel vector for singularity;
\item a characteristic-polynomial identity for Cayley--Hamilton.
\end{itemize}

For spectral constructions, distinguish stable data from labels.  Eigenvalue multisets, characteristic coefficients, and separated spectral projectors are stable.  Individual eigenvectors require a simple-eigenvalue gap and a normalization guard.  At multiplicity, output a subspace presentation or projector relation.

\subsection{Real and complex analysis}

A real number is represented by a primitive-recursive rational name \(r(p)\) and a regularity certificate.  A function on a compact rational box is represented by an evaluation algorithm plus a modulus of uniform continuity.  Differentiable or analytic objects carry explicit derivative, remainder, or coefficient bounds.

\paragraph{Arithmetic of names.}
Addition and multiplication compose demands and error budgets.  Inversion consumes apartness.  Maximum and absolute value are stable; exact sign is guarded or relational.

\paragraph{Integration.}
A Riemann or quadrature node computes a finite sum on a mesh selected from the modulus.  The local theorem bounds the discretization error.  The graph's backward pass divides the terminal error budget among evaluation and mesh errors.

\paragraph{Differentiation.}
A derivative name requires a supplied effective remainder estimate.  Mere pointwise differentiability does not produce a modulus.  Finite difference quotients are exact rational nodes; the mathematical producer supplies the error bound.

\paragraph{Power series and analytic branches.}
Coefficients and convergence radii must be explicit.  A local inverse or implicit branch requires a derivative/apartness guard and a quantitative contraction or Newton bound.  Exact branch coefficients from arbitrary Cauchy coefficient names are not a uniform operation; use exact algebraic input or an approximate relational branch.

\subsection{Complex integration and residues}

A contour is represented by finitely many rational polynomial or piecewise-linear segments with separation and smoothness data.  A finite quadrature computes the integral to requested precision.  Cauchy and residue arguments can be represented by:

\begin{itemize}
\item a finite triangulation or homotopy certificate avoiding poles;
\item local Taylor or Laurent coefficient bounds;
\item winding-number certificates on a finite grid;
\item explicit lower bounds for distance to singularities.
\end{itemize}

The topological content is finite and combinatorial; the analytic content is a modulus/remainder producer.  A statement that merely says the contour avoids a pole, without a positive distance certificate in a compact presentation, is not yet an executable interface.

\subsection{Optimization and convexity}

For a compact effective domain, a finite net and continuity modulus give an approximate minimizer by bounded argmin.  Convexity certificates can improve the bound but are not required by the generic compiler.

An exact minimizer is generally not computable from approximate data.  Appropriate outputs are:

\begin{itemize}
\item a point with objective within \(2^{-p}\) of the infimum;
\item an interval containing the infimum;
\item a finite cloud of approximate minimizers;
\item a unique minimizer name under a strong convexity or isolation modulus.
\end{itemize}

KKT or dual certificates are finite relational witnesses.  Their checker proves primal feasibility, dual feasibility, and a bounded gap.

\subsection{Elementary and projective geometry}

Coordinate geometry over rationals, signed integers, or explicit algebraic extensions is highly compatible with iPRA.  Incidence, collinearity, conic membership, cross ratios, determinants, and squared distances are polynomial identities.  Degenerate cases should be kept explicit rather than divided away.

A robust theorem statement has:

\begin{itemize}
\item exact point/line/conic codes;
\item Boolean nondegeneracy or incidence hypotheses;
\item a division-free polynomial conclusion where possible;
\item guarded rational operations when division is unavoidable.
\end{itemize}

Synthetic choice constructions, such as selecting an intersection or orientation, are represented by finite relations.  If the final theorem is a polynomial incidence identity, stable erasure removes the auxiliary choices.

\subsection{Trigonometry and geometric constants}

Transcendental functions are represented by rational names with explicit remainder bounds, not by exact function symbols.  Trigonometric identities can be proved by algebraic operations on simultaneous sine/cosine names or by finite Taylor certificates.  Division-based identities require denominator guards.

Angles themselves need not be canonical.  Unit vectors, dot and cross products, or half-angle parameterizations often give better finite interfaces.  On the sphere, use exact vector identities and introduce inverse trigonometric functions only for approximate measurement outputs.

Constants such as \(\pi\) are named by primitive-recursive rational sequences with proved error bounds, for example polygonal or arctangent constructions.  Equality to a geometric quantity is an observational statement between names.

\subsection{Combinatorics}

Finite combinatorial objects are naturally coded by bounded tables.  Counting functions, recurrences, generating-polynomial coefficients up to a finite degree, and congruences are primitive recursive.  Bijective proofs can be represented by mutually inverse finite transducers and checked identities.

Asymptotic formulas require explicit constants and remainder bounds.  For each request \(p\) and parameter range, a certificate checker may verify the finite inequalities.  A qualitative \(O(\cdot)\) statement with ineffective constants is not itself a PFIG contract.

Infinite generating functions are treated coefficientwise: fix a coefficient index, propagate a finite truncation demand, and prove the coefficient identity.  Universal generalization occurs after the finite computation.

\subsection{Elementary, algebraic, and analytic number theory}

Divisibility, gcd, primality certificates, modular arithmetic, finite products over primes below a bound, and recurrence sequences are primitive recursive.  Existential number-theoretic statements should expose witnesses or bounded searches.

For analytic number theory, use a certificate format containing:

\begin{itemize}
\item rational approximations to finite sums or integrals;
\item explicit truncation and tail bounds;
\item zero-free or nonvanishing margins;
\item finite interval verification data;
\item the exact order of quantifiers for constants and thresholds.
\end{itemize}

The generic verifier can be formalized in iPRA.  Producing sharp effective constants may remain deep theorem-specific work.

\subsection{Diophantine approximation}

Dirichlet-type theorems are finite pigeonhole arguments and fit bounded search.  Continued fractions are bounded loops with exact integer invariants.  A finite convergent list and approximation inequalities can be certified directly.

Statements that choose a best approximation from an infinite set must be reformulated at a finite denominator bound or equipped with a proof that a bounded search is complete.  Metric ``almost everywhere'' theorems require finite rational probability approximations and explicit measure-error bounds.

\subsection{Diophantine geometry}

A rational point is a normalized or unreduced integer tuple with a nonzero denominator and a polynomial equality certificate.  Bounded-height enumeration is primitive recursive.  Maps between varieties are guarded rational transducers, with exceptional loci recorded explicitly.

Finite descent has the PFIG form:

\begin{enumerate}
\item enumerate or receive a finite list of covering classes;
\item for each class, produce a relation of local/global representatives;
\item check compatibility and descent equations;
\item combine the surviving classes by finite viability;
\item obtain generators or point lists under an explicit completeness bound.
\end{enumerate}

The representative construction is not automatic.  A Selmer-class code is not by itself a rational point on the covering variety.  The framework exposes this as a producer node rather than treating it as a harmless decoding step.

\subsection{Probability}

The native objects should be finite rational laws or finite random variables.  Expectations, moments, convolutions, characteristic functions at rational arguments, and tail probabilities are finite sums.

For approximation of general distributions, require:

\begin{itemize}
\item a quantization map to a finite rational law;
\item a tail bound;
\item a moment or Lindeberg-type modulus;
\item anti-concentration data when a discontinuous CDF metric is used.
\end{itemize}

Central-limit statements then become explicit inequalities between finite distributions with a Berry--Esseen-type bound.  Weak convergence without a rate is not a regular-name interface.

\subsection{Topology}

Use finite simplicial or cubical approximations, carriers, chain complexes, and parity certificates.  Continuous maps are accessed only through finite observations and a modulus/carrier condition.  Homotopies are finite prism operators or bounded sequences of contiguous maps.

A Jordan, invariance-of-domain, or dimension theorem should identify the finite combinatorial invariant carrying the proof: winding parity, Sperner labeling, intersection number, or chain obstruction.  The continuum statement follows from the presentation's approximation theorem and separation data.

\subsection{Differential geometry and mathematical physics}

For coordinate calculations, metrics, connections, curvature tensors, and field equations are finite rational or symbolic expressions.  A checker can verify inverse-matrix, determinant, Bianchi, Maxwell, or Einstein identities exactly after clearing denominators.  Validity domains and nonzero determinants are guards.

Global geometric conclusions require atlas compatibility, coverage, and analytic existence data.  A symbolic local identity should not be advertised as a global manifold theorem without those interfaces.  Finite chart gluing and observational equality provide the correct bridge.

\subsection{Category-theoretic organization}

Category language is valuable as a metatheoretic organizer of contracts.  Exact contracts compose, stable observers are natural with respect to observational equivalence, and finite covers form diagram shapes.  For fixed finite presentations, Yoneda-style and finite-limit arguments can be compiled to bounded tables.

The categorical layer should not hide computational content.  Each hom-set used internally must be finite or represented by a fixed PR family with explicit evaluation.  Universal properties are certificate relations, not unbounded quantification over arbitrary programs.

\subsection{A standard theorem-packet template}

A repository theorem packet can be organized uniformly:

\begin{enumerate}
\item \textbf{Representation}: exact codes, observations, equivalence, restrictions.
\item \textbf{Demand}: precision/resource transformers and quantitative guards.
\item \textbf{Producer}: exact term, bounded loop, bounded search, relation, or imported certificate.
\item \textbf{Local object proofs}: base, step, reflection, extensionality, stability.
\item \textbf{Graph}: dependencies and terminal request.
\item \textbf{Compilation}: generated PR state, quantifier-free invariant, and terminal \(\Prf\ \PRAAx\) theorem.
\item \textbf{Audit}: theorem ledger distinguishing semantic, tested, object-level, and transferred claims.
\item \textbf{Barriers}: any missing modulus, gap, height, representative, or effective constant.
\end{enumerate}

This template is broad enough for the repository's current areas while remaining specific enough to prevent logical ambiguity.

\section{Resolution of reported difficulties and principled barriers}
\label{sec:difficulties}

This section maps the recurring repository difficulties to precise PFIG responses.  The goal is not to declare every deep theorem solved, but to identify which part is now generic, which part must be reformulated, and which part remains a theorem-specific producer.

\subsection{Ambient evidence mistaken for object proof}

\textbf{Problem.} A Lean theorem proves a semantic property of a finite computation, or Python tests pass, but the final theorem is stated only as a desired \(\Prf\ \PRAAx\) value.

\textbf{Resolution.} Every node contract stores object derivations for base, step, reflection, and terminal extraction.  PFIG soundness composes those derivations using the existing substitution and cut lemmas.  Semantic execution remains a reference specification and test oracle.

\textbf{Residual work.} Enter the theorem-specific arithmetic lemmas in object syntax.  The framework removes repeated control-flow proofs, not the local mathematics.

\subsection{Quantifiers in induction invariants}

\textbf{Problem.} Correctness is stated with bounded or unbounded quantifiers and therefore cannot be used as a QF induction matrix.

\textbf{Resolution.} Compile every finite quantifier to an \(\mathsf{All}\), \(\mathsf{Any}\), search, sum, product, or table fold.  Induct on a literal equality or Boolean state invariant.  Reflect the resulting Boolean only after the loop.

\textbf{Status.} Resolved as a general schema; a shared Lean Booleanization library remains to be completed.

\subsection{All precisions inside one invariant}

\textbf{Problem.} A proof attempts to maintain a statement about every future precision during a finite recursion.

\textbf{Resolution.} Fix the terminal request as a free parameter, run backward demand, instantiate a finite query table, prove the finite run, then universally generalize over the request.  Uniformity comes from the PR request parameter, not from an infinitary invariant.

\textbf{Status.} Resolved by the two-pass compiler.

\subsection{Universal primitive-recursive evaluator}

\textbf{Problem.} The construction treats arbitrary PR program codes as data and assumes a total PR evaluator.

\begin{barrier}[No total PR universal evaluator]
There is no total primitive-recursive function \(U(e,x)\) enumerating every unary primitive-recursive function by code \(e\).  Otherwise the diagonal function \(d(n)=U(n,n)+1\) would be primitive recursive and equal to some \(U(e,\cdot)\), giving \(d(e)=d(e)-1\).
\end{barrier}

\textbf{Resolution.} Statically specialize each fixed graph to fixed PR symbols, or verify a finite clocked trace.  The metatheoretic compiler may be general; the object theory receives specialized syntax.

\textbf{Status.} The stronger interface is impossible; the specialized interface is resolved.

\subsection{Non-unique witnesses and incoherent names}

\textbf{Problem.} Independent least choices at each precision jump among branches.

\textbf{Resolution.} Keep finite answer clouds, compute backward viability for fixed horizons, and either select a finite viable path or erase choices through a stable observer.  Construct a regular name only under explicit cloud-diameter and motion-tail bounds.

\textbf{Status.} Finite composition is resolved.  Infinite name extraction is conditional on quantitative shrinking data, as it must be.

\subsection{Root ordering and spectral labels}

\textbf{Problem.} Individual roots or eigenvectors are labeled across collisions.

\textbf{Resolution.} Return multisets, clusters, projective data, symmetric observables, or finite matchings.  Exact labels are derived under a separation gap using \cref{thm:branch-stability}.

\textbf{Status.} Reformulated in the unrestricted case; resolved under a gap.

\subsection{Exact sign, side, and axis selection}

\begin{barrier}[Discontinuous binary choice]
A total operation that, from an arbitrary Cauchy name of a real, decides whether it is nonnegative or nonpositive would imply a limited omniscience principle.  Similarly, a total exact rotation-axis selector at the identity must choose among an entire sphere of valid axes and cannot be continuous.
\end{barrier}

\textbf{Resolution.} Use a finite relation of admissible sides/axes, a projective or quaternion representation, a stable formula independent of the choice, or a positive separation guard.

\textbf{Status.} The unrestricted selector is intentionally not implemented.

\subsection{Partial operations from mere nonzero propositions}

\textbf{Problem.} A proof divides by a represented quantity after assuming only \(x\ne0\).

\textbf{Resolution.} Totalize the code function and require an apartness witness \(|x|\ge2^{-g}\).  Backward demand uses \(g\) to determine input precision.  If no margin is available, return a relation or retain the statement as noncomputational.

\textbf{Status.} Resolved under quantitative guards.

\subsection{Compactness and minimization}

\textbf{Problem.} ``Choose a finite subcover'' or ``take a minimizer'' appears without a modulus or finite net.

\textbf{Resolution.} Require a primitive-recursive net, coverage certificate, continuity modulus, and bounded argmin.  A unique exact minimizer additionally requires an isolation modulus.

\textbf{Status.} Resolved for effective compact presentations; the net/modulus producer is theorem-specific.

\subsection{Metastability treated as convergence}

\begin{barrier}[Metastability does not uniformly yield a Cauchy modulus]
Computable sequences can have primitive-recursive finite-window metastability information while lacking a computable tail modulus.  A uniform conversion would compute convergence rates for Specker-type sequences.
\end{barrier}

\textbf{Resolution.} Keep metastable output as a finite-window contract.  Upgrade to a regular name only with an explicit tail modulus, for example from shrinking viability.

\textbf{Status.} The distinction is enforced by separate interface types.

\subsection{Exact Puiseux or branch data from approximate coefficients}

\textbf{Problem.} A uniform algorithm is asked to output exact branch coefficients or exact factorization type from arbitrary Cauchy coefficient names.

\textbf{Resolution.} Use exact algebraic coefficient codes and splitting certificates, or produce approximate branch clouds.  Exact branch labels require discriminant/separation data.  Arbitrary Cauchy input supports only stable approximate observables.

\textbf{Status.} Reformulated; the unrestricted producer is obstructed by discontinuity and exact-zero detection.

\subsection{Finiteness confused with a height bound}

\begin{barrier}[Cardinality does not locate solutions]
A bound on the number of solutions of a decidable predicate does not provide a bound on where they occur.  If it did uniformly, one could encode a halting computation as a predicate with at most one solution and decide whether and where it halts.
\end{barrier}

\textbf{Resolution.} Separate the height producer from the bounded enumerator.  Use \cref{thm:height-to-list,thm:list-to-height} to move between explicit height bounds and complete lists.

\textbf{Status.} The verifier is resolved; general Faltings/Siegel height producers remain open or theorem-specific.

\subsection{Mordell--Weil and descent representatives}

\textbf{Problem.} A finite cohomological or Selmer-class list is treated as though it contained rational points on every covering variety.

\textbf{Resolution.} Make rational representatives and compatibility certificates explicit relational producer nodes.  Finite descent, group generation from a complete covering list, and certificate checking are generic after these nodes are supplied.

\textbf{Status.} Verifier and relational descent are resolved; uniform representative production remains an honest open obligation.

\subsection{Ineffective analytic constants}

\textbf{Problem.} An asymptotic theorem hides constants or thresholds not computable from the input.

\textbf{Resolution.} State a finite analytic certificate containing explicit rational constants, threshold, and error inequalities.  PFIG verifies the certificate and propagates precision.  A separate producer must generate the constants.

\textbf{Status.} Verification architecture resolved; effectiveness of the source theorem remains substantive mathematics.

\subsection{General probability spaces}

\textbf{Problem.} A proof invokes arbitrary measures, discontinuous CDFs, or weak convergence without quantitative approximation data.

\textbf{Resolution.} Work first with finite rational laws.  Approximation of a wider class requires quantization, tail, moment, and anti-concentration moduli.  The finite theorem and transfer error are separate nodes.

\textbf{Status.} Resolved for effective presentations.

\subsection{Continuous topology without numerical control}

\textbf{Problem.} A continuous compact space and map are used without a modulus, grid, carrier, or finite invariant.

\textbf{Resolution.} Represent each requested stage by a finite complex or grid, maps by carriers, and the theorem by parity, Sperner, winding, chain, or homology certificates.  Stable-observer arguments remove selection dependence.

\textbf{Status.} Representation and finite proof layer resolved; theorem-specific approximation and separation estimates remain.

\subsection{Code equality versus mathematical equality}

\textbf{Problem.} Equivalent representatives are required to be literally equal, causing false canonicity or impossible normalization obligations.

\textbf{Resolution.} Use observational setoids, finite back-and-forth for relation clouds, and explicit extensionality demands.  Quotients live in the semantics, not necessarily in the object syntax.

\textbf{Status.} Resolved as a general interface.

\subsection{Classical proof fragments}

\textbf{Problem.} A finite arithmetic argument uses classical case analysis, and it is unclear which conclusions survive intuitionistically.

\textbf{Resolution.} The negative translation of \cref{sec:negative-transfer} transfers classical PRA derivations whose assumptions and conclusions lie in the retained class.  Positive existential or disjunctive outputs must instead be made explicit by PFIG witnesses or relations.

\textbf{Status.} Human proof complete; Lean implementation remains an engineering target.

\subsection{A hierarchy of remaining obligations}

After PFIG compilation, unresolved work falls into four clearly separated classes.

\begin{enumerate}
\item \textbf{Kernel engineering}: implementing Booleanization, viability reflection, graph elaboration, and negative translation in Lean.
\item \textbf{Local arithmetic}: proving a finite inequality, identity, or certificate reflection theorem in object syntax.
\item \textbf{Quantitative producers}: deriving a modulus, gap, height bound, cover, representative, or explicit constant.
\item \textbf{Genuine barriers}: requests for discontinuous selectors, universal evaluators, or non-effective conversions that should be reformulated rather than pursued unchanged.
\end{enumerate}

This hierarchy is itself a useful research result.  It prevents effort from being spent on formalizing an interface that cannot exist, and it identifies which hard theorem would unlock an entire family of downstream iPRA results.

\section{Lean formalization architecture}
\label{sec:lean-plan}

This section gives a source-oriented plan against the checked \texttt{ConciseIPRA} API.  It is deliberately more precise than a list of theorem names: it fixes variable order, formula shapes, module dependencies, and the point at which each semantic construction becomes an actual object-iPRA derivation.

\subsection{Pinned upstream interface}

The intended imports are the repository modules providing:

\begin{itemize}
\item \texttt{Syn}, \texttt{Fm}, substitution, shifting, and de Bruijn utilities;
\item \texttt{Prf} and notation \texttt{Gamma |-[PRAAx] phi};
\item primitive-recursive syntax \texttt{PR} and terms \texttt{PTm};
\item the iPRA axiom family \texttt{PRAAx} and \texttt{PRA.qfInd};
\item the native transfer theorem \texttt{pra\_iff\_native}, used only after an ordinary derivation is complete.
\end{itemize}

The inspected repository examples import \path{ConciseIPRA.PR} and work in namespace \path{LeanProofs.ConciseIPRA}.  The packet's blueprint follows that convention.  Before integration, the exact commit should be pinned and every imported name checked against the current module tree.

\subsection{Module graph}

The proposed development is split into ten modules.

\begin{longtable}{@{}p{0.27\textwidth}p{0.65\textwidth}@{}}
\toprule
Module & Responsibility \\
\midrule
\endhead
\texttt{PFIG/Boolean.lean} & Boolean terms, range theorems, reflection of connectives, bounded \texttt{all}/\texttt{any}, least search, and QF formula Booleanization. \\
\path{PFIG/ExactContract.lean} & Exact term contracts, specialization, composition, and stable-observer combinators.  Most proofs adapt existing repository code. \\
\texttt{PFIG/ObjectLoop.lean} & Source-level wrapper around \texttt{primitiveRecursionInvariant} and \texttt{primitiveRecursionAt}; terminal extraction. \\
\texttt{PFIG/Demand.lean} & Coded request semilattices, monotone demand maps, fixed-DAG reverse pass, and object reflection of demand equations. \\
\texttt{PFIG/Relation.lean} & Finite relation codes, viability table layout, path checker, least viable path, and object reflection. \\
\texttt{PFIG/Shrinking.lean} & Coded rational metrics, viable-cloud diameter, edge motion, adjacent coherence, finite triangle sums, and regular-name extraction. \\
\texttt{PFIG/Guarded.lean} & Apartness guards, totalized operations, branch stabilization, observational setoids, and relational back-and-forth. \\
\texttt{PFIG/Graph.lean} & Fixed graph syntax or generated-stage syntax, topological schedule, two-pass state compiler, invariant assembly, and soundness theorem. \\
\path{PFIG/NegativeTransfer.lean} & Equality stability, stability of negative translations, proof translation, retained equivalence, and retained transfer. \\
\texttt{PFIG/Applications.lean} & Worked end-to-end examples and adapters for height/list, finite gluing, root clusters, and finite nets. \\
\bottomrule
\end{longtable}

The first complete milestone need not internalize the graph AST.  It may use generated theorem-specific stages and prove them with \texttt{ObjectLoop}.  \texttt{Graph.lean} can then abstract the repeated pattern.

\subsection{Exact contracts}

The exact interface should reuse the repository's type:

\begin{lstlisting}[language={},caption={Exact object contract}]
structure ExactContract {n : Nat}
    (P : Fm PTm n) (Q : Fm PTm (n + 1)) where
  witness   : PTm n
  certified : [P] |-[PRAAx] Q.subst (inst witness)
\end{lstlisting}

Required derived operations are:

\begin{lstlisting}[language={}]
theorem ExactContract.existsOutput :
  c : ExactContract P Q -> [P] |-[PRAAx] Fm.ex Q

def ExactContract.compose
    (c : ExactContract P Q)
    (d : ExactContract Q R) :
    ExactContract P (R.subst (up (inst c.witness)))

def ExactContract.specialize
    (c : ExactContract P Q) (sigma : Fin n -> PTm m) :
    ExactContract (P.subst sigma) (Q.subst (up sigma))
\end{lstlisting}

The inspected \texttt{ObjectComposition.lean} already gives these constructions.  The new module should import or relocate them rather than duplicate proofs.

\subsection{Loop contracts}

Use a formula \(P\) with free variables ordered as state, time, then parameters.  Define helper substitutions exactly once.

\begin{lstlisting}[language={}]
def instStateTime {k n : Nat}
    (P : Fm PTm (k + 2))
    (state time : PTm n)
    (params : Fin k -> PTm n) : Fm PTm n :=
  P.subst (Fin.cases state (Fin.cases time params))
\end{lstlisting}

A source-facing contract is:

\begin{lstlisting}[language={}]
structure LoopContract {k : Nat}
    (Gamma : List (Fm PTm k)) (Post : Fm PTm k) where
  init      : PR k
  step      : PR (k + 2)
  fuel      : PTm k
  invariant : Fm PTm (k + 2)
  qf        : invariant.QF
  base      : Gamma |-[PRAAx] baseFormula init invariant
  preserve  : invariant :: (Gamma.map Fm.shift).map Fm.shift
              |-[PRAAx] nextFormula step invariant
  finish    : Gamma |-[PRAAx]
              instStateTime invariant
                (runTerm init step fuel var) fuel var --> Post
\end{lstlisting}

Then:

\begin{lstlisting}[language={}]
theorem LoopContract.sound (c : LoopContract Gamma Post) :
  Gamma |-[PRAAx] Post :=
  Prf.impE c.finish
    (primitiveRecursionAt c.init c.step c.invariant
      c.qf Gamma c.base c.preserve c.fuel)
\end{lstlisting}

All substitution lemmas required by \texttt{primitiveRecursionAt} already appear in the inspected \texttt{ObjectIteration.lean}.  The integration task is to make that theorem a stable library module and remove copied \texttt{ObjectKernel} snapshots from individual packets.

\subsection{Boolean formulas}

There are two useful layers.

\paragraph{Term-level Boolean algebra.}
Fix PR symbols \texttt{bnot}, \texttt{band}, \texttt{bor}, \texttt{beq}, and bounded folds.  Prove object theorems such as:

\begin{lstlisting}[language={}]
theorem band_reflect (a b : PTm n) :
  [] |-[PRAAx]
    ((band a b = 1) <-> ((a = 1) /\ (b = 1)))

theorem bool_range (a : PTm n) :
  BoolTerm a -> [] |-[PRAAx] ((a = 0) \/ (a = 1))
\end{lstlisting}

The concrete formula constructors use the repository's notation, not the ASCII shown in the schematic listing.

\paragraph{Syntax-directed QF compiler.}
Define a metalevel function from a quantifier-free \(\Fm\) to a PR Boolean expression.  Because arbitrary formula syntax contains term symbols already interpreted by \texttt{PTm}, the compiler may return a \texttt{PTm n} plus a derivation record:

\begin{lstlisting}[language={}]
structure BoolReflection {n : Nat} (A : Fm PTm n) where
  code    : PTm n
  range   : [] |-[PRAAx] (code = 0 \/ code = 1)
  reflect : [] |-[PRAAx] Fm.iff (code = 1) A

def booleanizeQF {n : Nat} (A : Fm PTm n) (h : A.QF) :
  BoolReflection A
\end{lstlisting}

For a fully generic compiler, implication and disjunction cases use the already constructed decidability of subformulas.  A simpler first milestone provides reflection constructors and lets generated files assemble them.

\subsection{Bounded search}

Define a fixed PR state encoding a found bit and witness.  The theorem should expose both correctness and minimality:

\begin{lstlisting}[language={}]
structure SearchSpec {n : Nat} where
  bound : PTm n
  pred  : Fm PTm (n + 1)
  qf    : pred.QF
  bool  : BoolReflection pred

def leastSearch (s : SearchSpec) : PTm n

theorem leastSearch_sound
    (nonempty : Gamma |-[PRAAx] Fm.exBounded s.bound s.pred) :
  Gamma |-[PRAAx]
    s.pred.subst (inst (leastSearch s))

theorem leastSearch_minimal ...
\end{lstlisting}

If \texttt{Fm.exBounded} is not a library macro, define it as an ordinary existential conjoined with \(<\), and separately use the Booleanized bounded form for the loop invariant.

\subsection{Finite relations and table layout}

Keep object formulas independent of semantic finite types.  Use natural bounds and a fixed table encoding.

\begin{lstlisting}[language={}]
structure RelPresentation (n : Nat) where
  candBound : PTm (n + 1)       -- precision, params
  good      : Fm PTm (n + 2)    -- candidate, precision, params
  edge      : Fm PTm (n + 3)    -- next, current, precision, params
  goodQF    : good.QF
  edgeQF    : edge.QF
\end{lstlisting}

Document the variable order in comments and provide named substitution functions:

\begin{lstlisting}[language={}]
def RelPresentation.goodAt
    (R : RelPresentation n)
    (a p : PTm m) (xs : Fin n -> PTm m) : Fm PTm m

def RelPresentation.edgeAt
    (R : RelPresentation n)
    (b a p : PTm m) (xs : Fin n -> PTm m) : Fm PTm m
\end{lstlisting}

The viability table stores one bit per pair \((i,a)\) with \(i\le H\) and \(a<B_i\).  To avoid variable-length arrays in the object proof, use a rectangular bound
\[
 B^*(H,\vec{x})=\max_{i\le H}B_i(\vec{x})
\]
and mark out-of-range cells false.  Index \((i,a)\) by a primitive-recursive pairing function.  Reverse iteration can be implemented by the index \(H-i\).

Target theorems are:

\begin{lstlisting}[language={}]
theorem viableTable_reflect :
  Gamma |-[PRAAx]
    TableOK R H table -->
    (tableBit table i a = 1 <-> ViableFormula R H i a)

theorem leastViablePath_sound :
  Gamma |-[PRAAx] NonemptyViable R H -->
    PathOK R H (leastViablePath R H)
\end{lstlisting}

\texttt{ViableFormula} may contain bounded existentials and is used only after Boolean reflection; \texttt{TableOK} and the induction invariant remain QF.

\subsection{Shrinking metrics}

Represent nonnegative rationals by numerator/positive-denominator pairs or scaled integers.  Define a metric observation interface:

\begin{lstlisting}[language={}]
structure ObservedRelation (n : Nat) extends RelPresentation n where
  obs       : PR (n + 2)    -- candidate, precision, params -> rational code
  distLe    : Fm PTm (n + 3)
  diameter : PTm (n + 1)    -- precision, params -> rational code
  motion   : PTm (n + 1)
  diameterProof : ...       -- actual Prf PRAAx derivation
  motionProof   : ...
\end{lstlisting}

The exact arity of \texttt{distLe} depends on the rational-code API; the blueprint file keeps the theorem signature schematic.  The key object targets are:

\begin{lstlisting}[language={}]
theorem adjacent_coherence :
  Gamma |-[PRAAx]
    distLe (x p) (x (p+1)) (diameter p + motion p)

theorem regularName_of_tail :
  Gamma |-[PRAAx] TailBudget diameter motion --> RegularName x
\end{lstlisting}

The finite triangle-sum theorem is proved by \texttt{PRA.qfInd} on an equality/inequality code, not by invoking a semantic metric-space hierarchy.

\subsection{Observational setoids}

Use a structure containing formulas rather than Lean quotient types:

\begin{lstlisting}[language={}]
structure Presentation (n : Nat) where
  valid  : Fm PTm (n + 1)          -- code, params
  equiv  : Fm PTm (n + 3)          -- x, y, precision, params
  validQF : valid.QF
  equivQF : equiv.QF
  refl   : ...
  symm   : ...
  trans  : ...
  restrict : PR ...
  restrict_sound : ...
\end{lstlisting}

An operation record stores a code function, demand, guard, correctness, and extensionality derivation.  Relational operations replace a code function by a finite relation and extensionality by bounded back-and-forth.

\subsection{Demand semilattices}

Start with scalar requests \(\Nat\), where join is maximum.  After the scalar compiler works, generalize to coded tuples:

\begin{lstlisting}[language={}]
structure RequestCode where
  join : PR 2
  le   : PR 2
  refl : ...
  join_left  : ...
  join_right : ...
  join_least : ...
\end{lstlisting}

A fixed graph can initially carry plain Lean functions producing \texttt{PTm} demand expressions.  The object reverse pass is generated as one loop with a QF invariant.  Internalizing arbitrary graph syntax is not required for theorem-level soundness.

\subsection{Graph syntax and compilation}

After generated examples stabilize, define a typed metalevel graph.  One workable design separates numerical environments from proof obligations:

\begin{lstlisting}[language={}]
inductive NodeKind
  | exact | loop | search | relation | guard | observe | glue

structure Node where
  id        : Nat
  kind      : NodeKind
  deps      : List Nat
  request   : RequestSchema
  stateSize : Nat
  localCode : LocalProgram
  contract  : LocalContractSchema

structure Graph where
  nodes : List Node
  acyclic : AcyclicByIndex nodes
\end{lstlisting}

\texttt{LocalProgram} is not an arbitrary PR code evaluated internally.  It is metasyntax elaborated to a concrete \texttt{PR} term.  The compiler returns a package:

\begin{lstlisting}[language={}]
structure CompiledGraph (G : Graph) where
  demandInit demandStep demandFuel : ...
  execInit   execStep   execFuel   : ...
  demandInv execInv : Fm PTm ...
  demandQF : demandInv.QF
  execQF   : execInv.QF
  demandBase demandPreserve : Prf PRAAx ...
  execBase execPreserve     : Prf PRAAx ...
  extract : Prf PRAAx ...
\end{lstlisting}

\begin{lstlisting}[language={}]
theorem CompiledGraph.sound (C : CompiledGraph G) :
  InputContext G |-[PRAAx] TerminalFormula G := ...
\end{lstlisting}

The proof invokes the two loop soundness theorems and cuts \texttt{extract}.  Structural soundness of the compiler is an induction on \texttt{G.nodes}; acyclicity by index makes dependency lookup simple.

\subsection{Negative translation modules}

Use the existing definitions from \texttt{NegativeSyntax.lean}.  The target signatures are:

\begin{lstlisting}[language={}]
theorem equality_stable (s t : PTm n) :
  [] |-[PRAAx] ~~(s = t) --> (s = t)

theorem negative_stable (A : Fm PTm n) :
  [] |-[PRAAx] ~~(negative A) --> negative A

theorem translate_prf
    (h : Gamma |-[ClassicalPRAAx] A) :
  Gamma.map negative |-[PRAAx] negative A

theorem retained_iff
    (h : Retained A) :
  [] |-[PRAAx] Fm.iff (negative A) A

theorem retained_transfer
    (hGamma : forall B in Gamma, Retained B)
    (hA : Retained A) :
  (Gamma |-[ClassicalPRAAx] A) <-> (Gamma |-[PRAAx] A)
\end{lstlisting}

The actual syntax for double negation is \texttt{negFormula (negFormula A)}.  The \texttt{orE} and \texttt{exE} branches of \texttt{translate\_prf} should be implemented with helper lemmas:

\begin{lstlisting}[language={}]
theorem neg_or_elim
    (hOr : Gamma |- negative (A or B))
    (hA  : negative A :: Gamma |- C)
    (hB  : negative B :: Gamma |- C)
    (stableC : Gamma |- ~~C --> C) : Gamma |- C

theorem neg_ex_elim ...
\end{lstlisting}

This isolates the binder and context bookkeeping from the main induction.

\subsection{Worked formalization milestones}

The recommended order is proof-driven.

\begin{enumerate}
\item Re-export the existing exact contract, object iteration, and stable-observer files in one library without copied kernels.
\item Complete Boolean range/reflection for a small fixed set of PR Boolean symbols.
\item Formalize least bounded search and one example: choose the least divisor or least valid rational interval.
\item Formalize a tiny viability table with fixed level bounds, then generalize to variable bounds.
\item Prove shrinking coherence for exact rational observations and a geometric tail budget.
\item Generate and check one two-pass PFIG: bisection for \(y^3+y=x\), including backward precision demand and a regular root name.
\item Complete negative translation and retained transfer.
\item Add finite gluing and height/list adapters.
\item Only then internalize the general graph AST and prove compiler soundness.
\end{enumerate}

This order yields useful checked theorems early and keeps the generic compiler grounded in concrete examples.

\subsection{Axiom and status audit}

Every public theorem concluding \(\Prf\ \PRAAx\) should be followed in an audit file by \texttt{\#print axioms}.  Expected metatheoretic axioms inherited from Lean's quotient/propositional infrastructure should be recorded, and no new arithmetic axiom should be introduced.  Files containing \texttt{sorry}, \texttt{axiom}, or theorem propositions used as placeholders must be clearly marked blueprint-only.

The build should produce a machine-readable status file with, for each theorem:

\begin{itemize}
\item semantic Lean checked;
\item executable tests passed;
\item object formula defined;
\item \(\Prf\ \PRAAx\) derivation checked;
\item native transfer checked;
\item remaining producer assumptions.
\end{itemize}

This prevents a semantic theorem from being mistaken for completion of the object layer.

\subsection{Clear path to a complete formalization}

No new foundational principle is required.  The completion path consists of:

\begin{enumerate}
\item consolidating already existing object-level loop and contract code;
\item implementing finite Boolean reflection and table libraries;
\item proving viability and shrinking estimates by those libraries;
\item generating fixed PR state transitions and QF invariants;
\item composing local object certificates by substitution and cut;
\item completing the standard negative-translation induction;
\item applying the already checked ordinary-to-native equivalence at the end.
\end{enumerate}

The largest engineering risk is substitution bookkeeping for generated formulas, not proof-theoretic strength.  Named instantiation helpers, a fixed variable-order convention, and small regression examples should make that risk manageable.

\section{Conclusions and further research questions}
\label{sec:conclusion}

\subsection{Main conclusion}

The broad mathematical content of the repository does not require a separate foundational method for every subject.  A large common core is visible:

\begin{enumerate}
\item specialize a fixed construction to primitive-recursive syntax;
\item propagate a terminal request backward to a finite demand;
\item execute exact loops, bounded searches, and certificate checks forward;
\item keep non-unique data relational;
\item eliminate choices by finite viability or stable observation;
\item demand quantitative guards at discontinuities and partial operations;
\item encode the whole finite execution by a quantifier-free invariant;
\item return an actual \(\Prf\ \PRAAx\) derivation.
\end{enumerate}

The PFIG normal form turns this into one reusable backend theorem.  Its most important effect is methodological: it forces every hidden appeal to choice, compactness, exact zero testing, a convergence modulus, or a universal evaluator to appear as an explicit interface obligation.

\subsection{What has been resolved}

At the human mathematical level, this report gives complete arguments for:

\begin{itemize}
\item structural Booleanization of fixed bounded formulas;
\item least bounded search and finite choice;
\item proof-producing deterministic traces;
\item reverse demand and fixed-graph compilation;
\item finite relational viability and least viable paths;
\item coherent regular names from shrinking viable clouds;
\item exact and metric stable-observer elimination;
\item branch stabilization under a gap;
\item finite-cover gluing and representative-free global codes;
\item equivalence of explicit height bounds and complete finite lists;
\item negative translation and retained classical-to-intuitionistic transfer.
\end{itemize}

The report also explains why several unrestricted requests should be reformulated: exact root labels at collisions, exact axes at the identity, a universal PR evaluator, a Cauchy modulus from metastability alone, exact branch data from arbitrary approximate coefficients, and search cutoffs from mere finiteness.

\subsection{Immediate formalization questions}

\begin{enumerate}
\item What is the smallest object-iPRA Boolean library sufficient to reflect all bounded formulas used by current packets?  A syntax-directed compiler may be less work than a large hand-written arithmetic library.
\item Can the existing \texttt{ObjectIteration}, \texttt{ObjectComposition}, and \texttt{ObjectObservation} code be moved into a stable shared namespace without duplicating the upstream kernel?
\item Which state/table encoding minimizes de Bruijn substitution complexity: rectangular arrays, nested pairs, or bounded-sequence codes?
\item Should the first graph compiler emit specialized Lean files, or should a typed PFIG AST be internalized immediately?  The generated-file route appears lower risk.
\item Can proof terms for Boolean table reflection be compressed by using the native scalar calculus earlier, while preserving a readable ordinary-iPRA API?
\item How much of the negative-translation proof can be automated by generic natural-deduction lemmas independent of arithmetic?
\end{enumerate}

\subsection{Mathematical research questions}

\begin{enumerate}
\item Which classical theorems in the repository have retained conclusions after their existential content is converted to explicit PFIG witnesses?
\item Can root and spectral packets be uniformly reorganized around cluster relations and stable symmetric observables, with exact labels as a guarded corollary?
\item Which compactness arguments admit primitive-recursive finite-net producers with sharp enough moduli for practical formalization?
\item For which Diophantine families can explicit height producers be extracted in iPRA, thereby converting abstract finiteness packets into complete enumerators?
\item Can finite carrier relations and contiguity provide a uniform route from computational topology packets to object-iPRA homology theorems?
\item Is there a useful hierarchy of PFIG interfaces by logical strength---exact, bounded relational, metastable, regular-name, and guarded exact---that predicts which theorem transformations are possible?
\item Can proof mining systematically extract demand maps, tail budgets, and finite certificate formats from existing classical proofs?
\item Which error-budget schedules yield the smallest primitive-recursive state and proof terms for analysis applications?
\end{enumerate}

\subsection{Engineering research questions}

\begin{enumerate}
\item Can a small declarative contract language generate both Python reference code and Lean object formulas from one schema?
\item Can the compiler emit a theorem ledger automatically, marking every imported local certificate and every remaining producer assumption?
\item Can proof-carrying traces be shared across packets so that a checked determinant, polynomial, interval, or finite-net library becomes reusable?
\item Can native-calculus proof checking be used as a compact artifact format after ordinary-iPRA proofs are generated?
\item What is the best strategy for incremental proof checking when a graph contains large finite certificates: direct derivation trees, reflected Boolean computation, or a verified certificate VM specialized to a fixed instruction set?
\end{enumerate}

\subsection{Recommended repository consolidation}

A practical next repository step is to create a shared directory.
\begin{center}
\path{Logic/PrimitiveRecursiveArithmetic/PFIG/}
\end{center}
The directory should contain the modules from \cref{sec:lean-plan}.  Existing framework packets should then be treated as design documents and examples rather than copied kernels.  The theorem-packet README template should require:

\begin{itemize}
\item pinned upstream commit;
\item exact object formulas;
\item status of every \(\Prf\ \PRAAx\) target;
\item test commands and logs;
\item quantitative producer assumptions;
\item explicit barrier statements for intentionally unavailable selectors.
\end{itemize}

This would transform a collection of parallel proposals into one cumulative formalization program.

\subsection{Final perspective}

The right question is not whether ordinary mathematics can be copied verbatim into iPRA.  It is which finite observations, certificates, moduli, and stable relations constitute the proof's actual computational content.  Once those interfaces are made explicit, iPRA is strong enough to verify a surprisingly broad range of mathematics by primitive recursion and quantifier-free induction.  Where the interfaces cannot be supplied, the failure is informative: it identifies a discontinuity, a missing quantitative theorem, or a genuinely non-effective principle rather than an unexplained formalization impasse.

\appendix
\section{Theorem and status ledger}
\label{app:ledger}

The table below distinguishes human proofs in this report, executable finite evidence, inherited repository object theorems, and proposed Lean targets.  ``Human-complete'' does not mean kernel-checked.

\small
\begin{longtable}{@{}>{\raggedright\arraybackslash}p{0.08\textwidth}>{\raggedright\arraybackslash}p{0.52\textwidth}>{\raggedright\arraybackslash}p{0.29\textwidth}@{}}
\toprule
ID & Claim & Status \\
\midrule
\endhead
P1 & Structural Booleanization of every fixed bounded formula, with QF fold invariants. & Human-complete; Lean target. \\
P2 & Least bounded search, validity, and minimality. & Human-complete; Python-tested; Lean target. \\
P3 & Object trace invariant and specialization to a PR fuel. & Human proof; inherited object theorem in repository. \\
P4 & Trace-certificate reflection from local transition reflection. & Human-complete; Lean target. \\
P5 & Primitive-recursive reverse demand and minimality for fixed DAGs. & Human-complete; Python-tested; Lean target. \\
P6 & PFIG normal form as a query-specialized encoded loop. & Human-complete at construction level; generated Lean target. \\
P7 & Two-pass PFIG soundness from local object certificates. & Human-complete; depends on P1--P6 formalization. \\
P8 & Viability iff existence of a compatible finite path. & Human-complete; Python-tested; Lean semantic/object targets. \\
P9 & Least viable finite path and finite dependent choice. & Human-complete; Python-tested; Lean target. \\
P10 & Adjacent coherence from viable-cloud diameter and edge motion. & Human-complete; Python-tested. \\
P11 & Regular name from a primitive-recursive tail budget. & Human-complete; Lean rational-metric target. \\
P12 & Stable-observer erasure for independent finite certificates. & Human proof; inherited object combinator in repository. \\
P13 & Guarded branch stabilization under a positive gap. & Human-complete; Python-tested; Lean target. \\
P14 & Precision-contract composition. & Human-complete; exact case inherited; relational case target. \\
P15 & Certified finite-cover gluing. & Human-complete; Lean target. \\
P16 & Height bound implies complete finite list. & Human-complete; Python-tested; Lean target. \\
P17 & Complete finite list implies height bound. & Human-complete; Python-tested; Lean target. \\
P18 & Equality stability and stability of every negative translation. & Human-complete; Lean target. \\
P19 & Translation of every ClassicalPRAAx derivation to iPRA. & Human-complete; Lean target. \\
P20 & Retained equivalence and retained transfer. & Human-complete; Lean target. \\
B1 & No total primitive-recursive universal evaluator. & Diagonal barrier. \\
B2 & No uniform exact sign/axis selector at discontinuities. & Constructive discontinuity barrier. \\
B3 & Metastability alone does not yield a computable Cauchy modulus. & Computability barrier. \\
B4 & Finiteness/cardinality alone does not yield a search cutoff. & Halting-style barrier. \\
B5 & Exact arbitrary-Cauchy branch data is not uniformly available. & Discontinuity/exact-zero barrier. \\
O1 & Shared object-level Booleanization compiler in Lean. & Open engineering task. \\
O2 & Generic object viability table reflection. & Open engineering task. \\
O3 & Verified internal PFIG AST compiler. & Open engineering task; generated specialization is interim route. \\
O4 & General effective Faltings/Siegel height producer. & Open mathematical producer. \\
O5 & Uniform rational representatives for general Mordell--Weil descent inputs. & Open mathematical producer. \\
\bottomrule
\end{longtable}
\normalsize

\section{Notation and variable-order conventions}

For object formulas with parameters \(\vec{x}\), the formalization should use the following uniform de Bruijn order:

\begin{center}
\begin{tabular}{@{}ll@{}}
\toprule
Formula & Free-variable order \\
\midrule
Loop invariant & state, time, parameters \\
Good candidate & candidate, precision, parameters \\
Compatibility edge & next candidate, current candidate, precision, parameters \\
Metric comparison & first code, second code, tolerance, parameters \\
Guarded operation & guard data, input code, request, parameters \\
\bottomrule
\end{tabular}
\end{center}

Named instantiation helpers should be the only code that constructs the corresponding \texttt{Fin.cases} substitutions.  All later proofs should use those helpers and their substitution lemmas.  This convention is not mathematically essential, but it is likely to eliminate a large class of formalization errors.


\begin{thebibliography}{99}
\bibitem{repo-pra-readme}
V.~Reshetnikov, \emph{ProveIt: PrimitiveRecursiveArithmetic README}, repository snapshot \texttt{34816faca8aa605999290b41d26210753b5c3b38}, 2026.

\bibitem{repo-concise}
V.~Reshetnikov, \emph{Concise presentations of intuitionistic primitive recursive arithmetic}, ProveIt repository, 2026.

\bibitem{repo-deduction}
V.~Reshetnikov, \emph{ConciseIPRA/Deduction.lean}, object-language intuitionistic derivability and structural lemmas, ProveIt repository, 2026.

\bibitem{repo-local-global}
V.~Reshetnikov, \emph{ObjectIteration.lean}, proof-producing primitive recursion in object iPRA, ProveIt repository, 2026.

\bibitem{repo-coherent}
V.~Reshetnikov, \emph{ObjectComposition.lean}, exact term contracts and loop certificates, ProveIt repository, 2026.

\bibitem{repo-observable}
V.~Reshetnikov, \emph{ObjectObservation.lean}, object-iPRA stable-observer combinators, ProveIt repository, 2026.

\bibitem{repo-unified}
V.~Reshetnikov, \emph{Unified finitary-relational framework for mathematics in iPRA}, ProveIt research packet, 2026.

\bibitem{bishop-bridges}
E.~Bishop and D.~Bridges, \emph{Constructive Analysis}, Springer, 1985.

\bibitem{bridges-richman}
D.~Bridges and F.~Richman, \emph{Varieties of Constructive Mathematics}, Cambridge University Press, 1987.

\bibitem{troelstra-vandalen}
A.~S.~Troelstra and D.~van Dalen, \emph{Constructivism in Mathematics: An Introduction}, Vols.~I--II, North-Holland, 1988.

\bibitem{kohlenbach}
U.~Kohlenbach, \emph{Applied Proof Theory: Proof Interpretations and their Use in Mathematics}, Springer, 2008.

\bibitem{kleene}
S.~C.~Kleene, \emph{Introduction to Metamathematics}, North-Holland, 1952.
\end{thebibliography}

\end{document}
